% Genuinely multipartite bound entanglement of three qubits.
% Build: pdflatex fully_ppt_three_qubit_gme.tex (twice).
\documentclass[11pt]{article}
\usepackage[T1]{fontenc}
\usepackage[utf8]{inputenc}
\usepackage{lmodern}
\usepackage[a4paper,margin=23mm]{geometry}
\usepackage{amsmath,amssymb,amsthm,mathtools}
\usepackage{booktabs,array,enumitem}
\usepackage{xcolor,xurl}
\usepackage[colorlinks=true,linkcolor=blue!45!black,citecolor=blue!45!black,urlcolor=blue!45!black]{hyperref}
\newtheorem{theorem}{Theorem}
\newtheorem{lemma}[theorem]{Lemma}
\newtheorem{proposition}[theorem]{Proposition}
\newtheorem{corollary}[theorem]{Corollary}
\newtheorem{conjecture}[theorem]{Conjecture}
\theoremstyle{remark}
\newtheorem{remark}[theorem]{Remark}
\setlength{\emergencystretch}{2em}
\allowdisplaybreaks
\newcommand{\ket}[1]{\lvert#1\rangle}
\newcommand{\bra}[1]{\langle#1\rvert}
\newcommand{\tr}{\operatorname{tr}}
\DeclareMathOperator{\rank}{rank}
\newcommand{\dg}{^{\dagger}}
\newcommand{\cls}[1]{}
\hypersetup{pdfauthor={Antony Lewis},pdftitle={Genuinely multipartite bound entanglement of three qubits}}
\title{Genuinely multipartite bound entanglement\\of three qubits}
\newcommand{\Sussex}{Department of Physics \& Astronomy, University of Sussex, Brighton BN1 9QH, UK}
\author{Antony Lewis\\[2pt]{\small\normalfont \Sussex}}
\date{}
\begin{document}
\maketitle
\begin{abstract}
If a three-party state is positive under partial transposition (PPT) across each of its three bipartitions, no local
protocol can distil any pure entanglement from it: neither singlets between any two parties nor GHZ states. All of its
entanglement is bound. We show that such a state can nevertheless be genuinely multipartite entangled (GME), meaning
that it cannot be prepared by mixing states in each of which some party is unentangled from the rest. This already
happens for three qubits, the smallest system in which the question arises and a common experimental setting
for multipartite entanglement; to our knowledge no three-qubit example was known. An explicit mixture of five orthogonal
states (one GHZ-type, three $W$-type and one product state) is PPT across every cut and GME throughout an exact parameter interval, and its rank, five, is
the minimum possible. Genuine entanglement is proved by a witness acting only on the five-dimensional support, whose
nonnegativity on the biseparable states in that support reduces to a sum of squares on its curves of product vectors. Such
states are invisible to the PPT-mixture test, the standard semidefinite relaxation used to certify GME, yet they form
a set with nonempty interior among density matrices: an explicit full-rank example remains GME under $2.33\%$ white noise, while $2.336\%$
makes it biseparable. We also construct extreme examples with the smallest possible partial-transpose ranks
$(5,5,5,5)$. All our explicit states, including the full-rank one, lie in the mixed $W$ class: they are mixtures of pure
states with vanishing three-tangle, so their genuine entanglement needs no GHZ-type component. Full PPT therefore does not imply
biseparability, even for three qubits.
\end{abstract}

\section{Introduction}\label{sec:ppt3-intro}

Entanglement among three or more parties comes in two kinds. A state is \emph{biseparable} if it is a mixture of states
each separable across one of the bipartitions $A|BC$, $B|AC$, $C|AB$; the bipartition may differ from term to term. Otherwise it is \emph{genuinely multipartite entangled} (GME). Each term of a biseparable mixture is entangled at most within one pair of
parties, so GME is the multipartite property that experiments with several qubits usually aim to
certify~\cite{HHH09,GT09}. Showing that a mixed state is GME is harder than showing it entangled across
each cut: a state can be entangled across every bipartition and still be a mixture of states separable across
different ones. Biseparability is a property of the convex hull of three sets of separable states, not of each set
separately.

Entanglement can also be distillable or bound. Entanglement across a cut is \emph{distillable} if many copies can
be converted to singlets across that cut by local operations and classical communication (LOCC), and \emph{bound}
otherwise. A state that is positive under partial transposition (PPT) across a cut is undistillable across it~\cite{HHH98},
and PPT is also necessary for separability~\cite{Peres,HHH96}. For three qubits $A$, $B$, $C$ we call a state
\emph{fully PPT} when its partial transposes across $A|BC$, $B|AC$ and $C|AB$ are all positive. Every LOCC protocol among
the three parties is in particular LOCC across each cut, and such operations preserve PPT across that cut~\cite{HHH98,HHH09}.
Consequently no pure state entangled across any cut, neither singlets between two parties nor GHZ states, can be
distilled from a fully-PPT state, even asymptotically.
Whatever entanglement a fully-PPT state has is bound.

This paper asks whether bound entanglement of this extreme kind can still be genuinely multipartite, and answers yes
for three qubits. A fully-PPT GME state carries entanglement that no single bipartition accounts for and from which no
pure entanglement can be distilled: it is present only collectively. The question also has a practical side. The standard
way to certify GME in experiments and numerics relaxes biseparability to \emph{PPT mixtures}, mixtures of states each PPT
across one cut, which gives a semidefinite program~\cite{JMG,NMG}. A fully-PPT state is already a one-term PPT mixture,
so its genuine entanglement is invisible to this test and to any test built from PPT conditions across cuts and their
mixtures. Novo, Moroder and G\"uhne conjectured that for three qubits PPT mixtures and biseparable states coincide~\cite{NMG};
Ha and Kye refuted this with a state PPT across one chosen cut~\cite{HK}. A fully-PPT GME state is the strongest form of
the counterexample.

Three qubits are where such a state is least expected. Each cut is $2\times4$, just beyond the $2\times2$ and $2\times3$
systems where PPT and separability coincide~\cite{HHH96}, and low rank forces separability: a fully-PPT three-qubit
state of rank at most four is separable across a cut, hence biseparable~\cite{KCKL,HLVC,KL,ABLS}. A fully-PPT GME state therefore
has rank at least five. Examples are known for more qubits and for higher local dimension (Section~\ref{sec:ppt3-prior}).
For three qubits, earlier constructions and numerical states were entangled but not shown to be GME, or were PPT across
one cut only~\cite{HK,KRB,GIM}, and Singh and Banik noted that no explicit example is known and derived necessary
rank conditions for one~\cite{SB}.

Our method follows the physics of the convex hull. Every pure state in a decomposition of $\rho$ lies in its support,
so a biseparable decomposition of a low-rank state can use only those vectors of the support that are products across some
cut. We choose a five-dimensional support $R$, spanned by five orthogonal vectors (one GHZ-type, three $W$-type and $\ket{111}$), in which the cut-product
vectors form one explicit curve per cut, parametrised by a single complex number and related by cyclic permutation of the
qubits. A Hermitian form that is nonnegative on these three curves is nonnegative on every biseparable state supported on
$R$; its negative expectation in our state then proves GME. For the form we use, nonnegativity on the curves is an explicit
factorisation into squares. Full PPT is a separate and elementary step, since cyclic and parity symmetry reduce each
partial transpose to two $4\times4$ blocks. Robustness to noise needs a test stronger than PPT; we use exact rational
certificates from the second level of the symmetric-extension hierarchy, together with explicit biseparable decompositions
that bound the noise thresholds from above.

Our main results are the following.
\begin{enumerate}[label=(\arabic*),itemsep=0pt,topsep=2pt]
\item \emph{Minimal-rank family} (Theorem~\ref{thm:ppt3-family}). A mixture $\rho_t$ of five orthogonal
vectors (GHZ-type, $W$-type and $\ket{111}$) has rank five and is fully PPT and GME for $3.257\lesssim t\lesssim3.389$ (exact endpoints in the
theorem).
\item \emph{A witness that sees only the support} (Proposition~\ref{prop:ppt3-witness}). One diagonal Hermitian form,
nonnegative on the product curves by a square factorisation, excludes biseparable mixtures across cuts, with no
semidefinite-program solver.
\item \emph{Extreme states} (Sections~\ref{sec:ppt3-family} and~\ref{sec:ppt3-surface}). Extreme fully-PPT GME states have
partial-transpose rank profiles $(5,6,6,6)$ and, the smallest possible, $(5,5,5,5)$; the latter lie in an analytic family
that is fully PPT and GME on the whole window $1/\sqrt3<\tau<1$ (one parameter modulo local filters). Every state of the family is
extreme in the fully-PPT set, and near every member the extreme states of this profile form a $27$-dimensional manifold with
nine local-filter moduli.
\item \emph{One support, both phases} (Section~\ref{sec:ppt3-coexistence}). One rank-five support carries both
biseparable and GME fully-PPT states, separated by an exact transition inside the strictly PPT region, so the support alone
does not decide genuine entanglement.
\item \emph{Robustness} (Section~\ref{sec:ppt3-noise}). Fully-PPT GME states have positive volume. An explicit full-rank
fully-PPT state remains GME under $2.33\%$ white noise, $1\%$ independent depolarisation per qubit, or a trace-distance
perturbation of $7.3\times10^{-3}$; its white-noise threshold is known to within half a per cent in relative terms. A
surface state of the $(5,5,5,5)$ family tolerates about $0.43\%$.
\item \emph{Structure} (Sections~\ref{sec:ppt3-chirality} and~\ref{sec:ppt3-relaxations}). No fully-PPT GME state is invariant
under the exchange of two qubits, so none is permutation invariant~\cite{NMG}, and our family is not locally equivalent
to its complex conjugate, not even by invertible local filters. Every explicit state we construct, including both extreme states
and the full-rank state $\tilde\rho$, belongs to the mixed $W$ class (Theorem~\ref{thm:ppt3-w-class}), as conjectured for PPT entangled three-qubit states by Ac\'in et al.~\cite{ABLS}. For every rank-five fully-PPT GME state the product vectors of its support form three rational
quartic curves, so supported witnesses of the kind used here detect all of them
(Theorem~\ref{thm:ppt3-universal-moments}).
\end{enumerate}

Section~\ref{sec:ppt3-relaxations} places the results relative to PPT mixtures and earlier work. The main
text gives the proof strategy; the computations are in Appendices~\ref{sec:ppt3-app-family}, \ref{sec:ppt3-app-surface},
\ref{sec:ppt3-app-coexistence} and~\ref{sec:ppt3-app-noise}, in the order of the sections they serve. Appendices~\ref{sec:ppt3-app-moment}
and~\ref{sec:ppt3-kernel} prove the product-kernel lemma and the moment cone of Section~\ref{sec:ppt3-relaxations} and the kernel-penalty
lift of the supported witness, and Appendix~\ref{sec:ppt3-verifiers} lists the proof dependencies. All theorems,
propositions, lemmas and corollaries are proved, analytically or with the exact computations named next; numerical
observations appear only in remarks and are labelled as such. Exact rational or algebraic
computations are used only for the extremality of $\sigma_6$ and the independent extremality check on the $(5,5,5,5)$ surface, the moduli count, the phase-coexistence minors and the noise certificates. The verification scripts with their
captured output are in a companion repository~\cite{Companion}, together with a Lean~4 formalisation of the main analytic
results and of the extremality of one state of the surface; the noise certificates, the moduli count, the extremality of
$\sigma_6$ and of the whole surface, and the minimality of rank five are checked by the exact scripts or cited, not formalised.

\section{The rank-five family}\label{sec:ppt3-family}

We first separate two tasks: enforce PPT by symmetry, then exclude biseparability using only product vectors in the
state's support. Appendix~\ref{sec:ppt3-app-family} supplies the coordinate proofs.

Let $\omega=e^{2\pi i/3}$. In the computational basis of $(\mathbb C^2)^{\otimes3}$ define the five
\emph{unnormalised} vectors
\begin{align}
u&=\ket{000}-\ket{011}-\ket{101}-\ket{110}, & v&=\ket{111}, & w&=\ket{001}+\ket{010}+\ket{100},\notag\\
e&=\ket{011}+\omega\ket{101}+\omega^2\ket{110}, & o&=\ket{100}+\omega\ket{010}+\omega^2\ket{001}. &&\label{eq:ppt3-vectors}
\end{align}
Here $u/2$ is a local-Clifford rotation of the GHZ state and $w,e,o$ are locally related $W$ states.
Global parity is $Z^{\otimes3}$, with $Z=\operatorname{diag}(1,-1)$; it has eigenvalue $+1$ on $u,e$ (an even number of
ones) and $-1$ on $v,w,o$. The cyclic permutation $\Pi$ of the qubits, $\Pi\ket{abc}=\ket{cab}$, fixes $u,v,w$ and
multiplies $e$ and $o$ by $\omega^2$. Together, $Z^{\otimes3}$ and $\Pi$ split $(\mathbb C^2)^{\otimes3}$ into two
two-dimensional $\Pi$-invariant sectors, $\operatorname{span}\{\ket{000},\ket{011}+\ket{101}+\ket{110}\}$ and
$\operatorname{span}\{v,w\}$, and four one-dimensional sectors spanned by $e,\bar e$ (even) and $o,\bar o$ (odd).
Exchanging two qubits swaps $e\leftrightarrow\bar e$ and $o\leftrightarrow\bar o$, and a PPT state invariant under such an
exchange is separable across a cut (Corollary~\ref{cor:ppt3-symmetrisation}). The support below therefore keeps $u$ from
the first sector, all of the second, and only one orientation $e,o$ of the one-dimensional sectors.
Put $B=[u,v,w,e,o]$, an $8\times5$ matrix with $B\dg B=\operatorname{diag}(4,1,3,3,3)$, and $R=\operatorname{ran}B$.
For $t>0$ set
\begin{equation}\label{eq:ppt3-rho}
\rho_t=\frac{t\,uu\dg+3t\,vv\dg+3\,ww\dg+3\,ee\dg+oo\dg}{7(t+3)} .
\end{equation}
For a three-qubit operator $X$ write $X^{T_c}$ for its partial transpose on qubit $c\in\{A,B,C\}$; the
\emph{rank profile} of $X$ is $(\rank X,\rank X^{T_A},\rank X^{T_B},\rank X^{T_C})$. For a vector $x$, $\|x\|$ is the
Euclidean norm; for an operator $X$, $\|X\|_1$ is the trace norm, $\|X\|_2=(\tr X\dg X)^{1/2}$ the Hilbert--Schmidt norm
and $\|X\|_\infty$ the operator norm (largest singular value).

\begin{theorem}[the rank-five fully-PPT GME family]\label{thm:ppt3-family}
\cls{proved; CF-1, CF-2, CF-3, IN-1, IN-2, IN-3, IN-4, IN-5, IN-6, IN-7}
Every $\rho_t$ has rank five. It is PPT at all three single-qubit cuts exactly
when
\begin{equation}\label{eq:ppt3-interval}
t_-\le t\le t_+,\qquad t_-=\frac{15+9\sqrt{17}}{16}\approx3.257,\qquad t_+=\frac{2(\sqrt{37}-1)}{3}\approx3.389,
\end{equation}
and every state in this interval is genuinely multipartite entangled (GME). In the interior the rank profile is
$(5,8,8,8)$, and at the endpoints it is $(5,7,7,7)$.
Rank five is minimal among fully-PPT GME three-qubit states.
\end{theorem}

For minimality, any fully-PPT three-qubit state of rank at most
four is separable across a $2\times4$ cut by the low-rank PPT
theorem~\cite{KCKL,HLVC}, and is therefore biseparable.

\emph{Idea of the PPT proof.} Cyclic invariance makes the three partial transposes unitarily equivalent, and global parity
splits $\rho_t^{T_A}$ into two $4\times4$ blocks (the partial transpose commutes with the real diagonal $Z^{\otimes3}$, whose eigenspaces
are spanned by $\ket{000},\ket{011},\ket{101},\ket{110}$ and by $\ket{001},\ket{010},\ket{100},\ket{111}$). The even-block
determinant forces $t\ge t_-$. In that range the odd-block determinant requires $t\le t_+$, and the first three leading
minors of both blocks are positive on the resulting interval. Sylvester's criterion then gives PPT. At an endpoint one block drops
to rank three, which gives the profiles. The minors are in Appendix~\ref{sec:ppt3-app-family}; genuine entanglement is proved
below by a witness on the range.

\subsection{Cut-product rays and a supported witness}\label{sec:ppt3-witness}

A biseparable decomposition of a state supported on $R$ uses pure vectors in $R$ that are products across some
cut. A Hermitian form nonnegative on all these vectors can therefore exclude biseparability on $R$.
For $s>0$ put
$u_s=\ket{000}-s(\ket{011}+\ket{101}+\ket{110})$, $B_s=[u_s,v,w,e,o]$ and $R_s=\operatorname{ran}B_s$, so that $u_1=u$,
$B_1=B$, $R_1=R$. An operator supported on $R_s$ is written $B_sXB_s\dg$ with a $5\times5$ coefficient matrix $X$;
subscripts identify the corresponding state. A pure cut-product vector $B_sq$ has
projector $B_sqq\dg B_s\dg$, so its coefficient matrix is $qq\dg$ and $\tr(H\,qq\dg)=q\dg Hq$ with no metric factor
$B_s\dg B_s$.

\begin{proposition}[supported witnesses]\label{prop:ppt3-witness}
\cls{proved analytically; CF-2, CF-3, IN-2, IN-3, IN-4, SW-1, SW-2}
Let $s>0$.
\begin{enumerate}
\item[(a)] The $A|BC$-product rays in $R_s$ are the rays of $B_sq_s(z)$, $z\in\mathbb C$, with
\begin{equation}\label{eq:ppt3-curve}
q_s(z)=\bigl(2z,\ -z^2(z^2+3s),\ z^2-s,\ -z(z^2+s),\ z^2+s\bigr)^T,
\end{equation}
together with the ray of $v$ (the point $z=\infty$). The other two cuts are given by the cyclic images.
\item[(b)] Let $H_s=\operatorname{diag}(-12s^3,\ 1,\ 11s^2,\ -6s,\ 14s^2)$ in the basis $(u_s,v,w,e,o)$. For
$z=x+iy$ and $r=|z|^2$,
\begin{equation}\label{eq:ppt3-sos}
q_s(z)\dg H_sq_s(z)=(r-s)^2\bigl[(r-5s)^2+12s\,x^2\bigr]\ \ge0 .
\end{equation}
At $s=1$ this is the witness $H_5:=H_1$ used for the first family,
\begin{equation}\label{eq:ppt3-h5}
H_5=\operatorname{diag}(-12,1,11,-6,14),
\qquad
q_1(z)\dg H_5q_1(z)=(r-1)^2\bigl[(r-5)^2+12x^2\bigr].
\end{equation}
Since $H_s$ is diagonal, it is invariant under the cyclic phases acting on $(u_s,v,w,e,o)$. Consequently
$\tr(H_sX)\ge0$ whenever $B_sXB_s\dg$ is biseparable, including mixtures across different cuts;
no symmetry hypothesis on $X$ is needed.
\end{enumerate}
\end{proposition}

\emph{Idea of the proof.} (a) A vector of $R_s$ is a product across $A|BC$ with first-qubit ray $(1,z)^T$ exactly when
its two $BC$ rows are proportional with ratio $z$: four linear equations in the five coordinates, with a one-dimensional
solution $q_s(z)$ for each finite $z$. (b) Since $H_s$ is diagonal, only the squared moduli of the entries of $q_s(z)$
enter, and the sum collapses to a perfect-square factorisation; a rescaling removes $s$. Diagonal forms commute with the
cyclic phases, so the other cuts give the same values, and every positive decomposition of an operator supported on $R_s$
uses vectors in $R_s$, so mixtures across cuts are controlled. The computations are in Appendix~\ref{sec:ppt3-app-family}.

\begin{proof}[Proof of the genuine entanglement in Theorem~\ref{thm:ppt3-family}]
Write $\rho_t=BX_tB\dg$. Then
\begin{equation}\label{eq:ppt3-h5-value}
\tr(H_5X_t)=\frac{29-9t}{7(t+3)}<0\qquad\text{whenever }t>29/9,
\end{equation}
and $t_->29/9$. Proposition~\ref{prop:ppt3-witness}(b) therefore proves GME throughout \eqref{eq:ppt3-interval}.
\end{proof}

For a concrete member, take $t=10/3$. Then the normalised versions of the orthogonal vectors
$(u,v,w,e,o)$ occur with probabilities
$(40,30,27,27,9)/133$, all leading minors in
\eqref{eq:ppt3-minors} are positive, and
$\tr(H_5X_{10/3})=-3/133$. Thus this fully-PPT GME state is certified by hand.
Appendix~\ref{sec:ppt3-kernel} lifts the supported witness to a global witness with an explicit noise radius around
this state (Remark~\ref{rem:ppt3-hand-radius}), again without a semidefinite-program solver.

\subsection{Reflection symmetry and local chirality}\label{sec:ppt3-chirality}

The states select one cyclic phase orientation of the $W$ vectors $e,o$; exchanging two qubits reverses it.
The following obstruction explains why this asymmetry matters.

\begin{corollary}[reflection symmetry destroys the construction]\label{cor:ppt3-symmetrisation}
\cls{proved; recalled swap-PPT reduction~\cite{NMG}}
\begin{enumerate}[label=(\alph*)]
\item A three-qubit state that is PPT at $A|BC$ and commutes with the swap of $B$ and $C$ is separable across $A|BC$.
Hence no fully-PPT GME state is invariant under the exchange of two qubits.
\item For every fully-PPT three-qubit state $\rho$, the permutation average
$\mathcal S(\rho)=\frac16\sum_{\pi\in S_3}U_\pi\rho U_\pi^\dagger$ is separable across each of the three cuts, so
\begin{equation}\label{eq:ppt3-asymmetry-upper}
 D(\rho,\mathrm{BS})\le\tfrac12\|\rho-\mathcal S(\rho)\|_1 ,
\end{equation}
where $D(\rho,\sigma)=\frac12\|\rho-\sigma\|_1$ is the trace distance, $D(\rho,\mathrm{BS})=\inf_{\sigma\in\mathrm{BS}}D(\rho,\sigma)$ and
$\mathrm{BS}$ is the set of biseparable states.
If $\rho$ is cyclically invariant and $S$ is any transposition of two qubits, the bound is
$\frac14\|\rho-S\rho S^\dagger\|_1$.
\item For any three-qubit $\rho$ that is PPT at $A|BC$, pinching with the triplet and singlet projectors $P_3$, $P_1$ of $BC$
gives an explicit $A|BC$-separable state, and
\begin{equation}\label{eq:ppt3-pinch}
D(\rho,\mathrm{BS})\le\|P_3\rho P_1\|_1 .
\end{equation}
\end{enumerate}
These are upper bounds on the distance to $\mathrm{BS}$, not entanglement tests. Coherence between singlet and triplet,
and the breaking of every pair-exchange symmetry, are therefore necessary for fully-PPT GME.
\end{corollary}
\begin{proof}
(a) The swap-invariant state is block diagonal in the triplet ($\dim3$) and singlet ($\dim1$) subspaces of $BC$, and its
partial transpose on $A$ preserves these blocks. Each block is PPT on $2\times3$ or $2\times1$, where PPT implies
separability; this is the mechanism of Novo, Moroder and G\"uhne~\cite[Sec.~III.2]{NMG}. (b) $\mathcal S(\rho)$ is fully
PPT and invariant under every transposition, so (a) applies at each cut; $\mathcal S(\rho)\in\mathrm{BS}$ gives
\eqref{eq:ppt3-asymmetry-upper}. For cyclic $\rho$, $\mathcal S(\rho)=(\rho+S\rho S^\dagger)/2$. (c) The pinched state
$P_3\rho P_3+P_1\rho P_1$ is PPT at $A|BC$ (the partial transpose on $A$ commutes with the pinching on $BC$) and block diagonal as
in (a), hence separable; its trace distance from $\rho$ is at most $\|P_3\rho P_1\|_1$. Since the swap is
$S_{BC}=P_3-P_1$, the pinched state equals $(\rho+S_{BC}\rho S_{BC})/2$ and $\|P_3\rho P_1\|_1=\frac14\|\rho-S_{BC}\rho S_{BC}\|_1$,
so (c) extends the cyclic bound of (b) to every state that is PPT at $A|BC$.
\end{proof}

The obstruction concerns reflection alone. The real vectors $u_s,v,w$ and the pair $e,o$ of \eqref{eq:ppt3-vectors} are invariant under the
antiunitary map $K\circ S_{BC}$, the swap of $B$ and $C$ followed by complex conjugation in the computational basis
(for example $S_{BC}e=\ket{011}+\omega^2\ket{101}+\omega\ket{110}=\bar e$). Hence every state with real coefficients in
these vectors, including $\rho_t$ and all states of Section~\ref{sec:ppt3-surface}, is invariant under this combined
reflection and conjugation. The GME states among them are invariant under neither separately: not under the reflection by
part~(a), and not under conjugation, which exchanges $e$ and $o$ with $\bar e$ and $\bar o$, outside their supports.

\begin{corollary}[intrinsic local chirality]\label{cor:ppt3-chirality}
\cls{proved; CF-4, SC-1--SC-5}
\begin{enumerate}[label=(\alph*)]
\item Throughout \eqref{eq:ppt3-interval}, $\rho_t$ and its complex conjugate are not related by local unitaries. The same
holds for physical spin time reversal, which differs from complex conjugation by a local unitary.
\item For every $s>0$, no invertible local filter $F_A\otimes F_B\otimes F_C$ maps $R_s$ onto its complex conjugate. Hence no state
with range $R_s$, in particular no $\rho_t$, $\sigma_6$, $\rho_{\tau,a}$ or state of the phase-coexistence line, is related to its complex
conjugate or its spin time reversal by an invertible local filter. The stabiliser of $R_s$ in $\mathrm{PGL}(2,\mathbb C)^{\times3}$ is
isomorphic to $S_3$, acting by the same M\"obius map on the three qubits.
\end{enumerate}
\end{corollary}
\emph{Idea of the proof.} (a) The eigenrays of $u$ and $v=\ket{111}$ are nondegenerate, so a local unitary relating $\rho_t$
to $\overline{\rho_t}$ must fix both, leaving only $I$ and $Z^{\otimes3}$; both fix $\rho_t$, which is not real
($\langle011|\rho_t|101\rangle\propto t+3\omega^2$). (b) $R_s$ contains exactly three fully product rays, $[\ket1^{\otimes3}]$ and
$[(\ket0\pm i\sqrt s\ket1)^{\otimes3}]$, and a filter must permute them; a M\"obius map is fixed by its action on three points, so the filter is
$F_0^{\otimes3}$ up to a scalar, which commutes with the cyclic shift. The shift has eigenvalue $\omega^2$ on $e,o\in R_s$ but not on $\overline{R_s}$.
See Appendix~\ref{sec:ppt3-app-family}.

\subsection{\texorpdfstring{An extreme point of profile $(5,6,6,6)$}{An extreme point of profile (5,6,6,6)}}\label{sec:ppt3-extreme6}

Adjusting the weight of $v$ makes both partial-transpose parity blocks singular and yields an extreme fully-PPT state.

\begin{corollary}[an extreme fully-PPT GME state of profile $(5,6,6,6)$]\label{cor:ppt3-extreme6}
\cls{exact-symbolic certificate, computer-assisted; the GME part is proved analytically; CF-6}
Set $a=t_-$ and $b=8(9-a)/5$ ($a\approx3.257$, $b\approx9.189$). Then
\begin{equation}\label{eq:ppt3-sigma6}
\sigma_6=\frac{a\,uu\dg+b\,vv\dg+3\,ww\dg+3\,ee\dg+oo\dg}{4a+b+21}
\end{equation}
is an extreme point of the set of fully-PPT three-qubit states and is GME. Its coefficient matrix is
$X_6=\operatorname{diag}(a,b,3,3,1)/(4a+b+21)$.
\end{corollary}
\emph{Idea of the proof.} The choice $a=t_-$ makes the even partial-transpose block singular, and that of $b$ makes the odd
block singular, each with a positive $3\times3$ leading minor, so every partial transpose has rank six. GME follows from the
witness $H_5$. By the supported-face criterion of~\cite{GIM}, extremality means that the only Hermitian $D$ supported on
$\operatorname{ran}\sigma_6$ whose partial transposes annihilate the kernels of those of $\sigma_6$ is a multiple of
$\sigma_6$; an exact rank computation on the $25$-dimensional space of coefficient matrices shows this
(Appendix~\ref{sec:ppt3-app-family}).

\section{\texorpdfstring{Profile $(5,5,5,5)$}{Profile (5,5,5,5)}: local-filter orbits}\label{sec:ppt3-surface}

How little partial-transpose rank can a fully-PPT GME state have? We deform the GHZ-type vector and tune the weights to
make both partial-transpose parity blocks singular. One of the two remaining parameters is an invertible local filter,
so a single parameter controls PPT and GME. Appendix~\ref{sec:ppt3-app-surface} gives the proofs.

The first family has partial-transpose ranks seven or eight. A fully-PPT state of rank at most four is
separable~\cite{KCKL,HLVC,KL,ABLS}, and the same theorem bounds each partial-transpose rank from below: if $\rank\rho^{T_c}\le4$, then $\rho^{T_c}$, and hence $\rho$, is separable across that cut~\cite{KCKL,HLVC,KL}.
Thus $(5,5,5,5)$ is the smallest possible profile for fully-PPT GME. The construction below attains it, including at
extreme points of the fully-PPT set.

For real $\alpha>0$, $0<\beta<\alpha$ and $\gamma>0$ put
$h=\sqrt{\alpha^2+\beta^2}$, $\kappa=\alpha^2-\beta^2$, $s=\alpha/\gamma$ and
\begin{equation}\label{eq:ppt3-weights}
O=\frac \kappa\beta,\qquad E=\frac{\kappa(h+\beta)^2}{\alpha^2\gamma},\qquad G=\frac{\alpha(E-\alpha s)}\beta,\qquad
V=\frac{\alpha^2E^2-2\alpha s\kappa E+(\alpha s)^2\kappa}{\beta\kappa},
\end{equation}
\begin{equation}\label{eq:ppt3-rho-surface}
\rho(\alpha,\beta,\gamma)=\gamma\,u_su_s\dg+\beta\,ww\dg+V\,vv\dg+G\,(wv\dg+vw\dg)+E\,ee\dg+O\,oo\dg,
\end{equation}
with $u_s,v,w,e,o$ as in Section~\ref{sec:ppt3-family}. The coefficients are chosen
to make both partial-transpose blocks singular. Where they come from (at $\alpha=\gamma=1$, so $\beta=\tau$): the odd block of $\rho^{T_A}$ has positive
$2\times2$ leading minor $3\tau O$, so it has rank two only if its $3\times3$ leading minor $3O\,(\tau(\tau+O)-1)$ vanishes, which gives $O=\kappa/\tau$; the last
row and column must then lie in the span of the others, which gives $G=(E-1)/\tau$ and $V=E^2/O+\tau G^2=(E^2-2\kappa E+\kappa)/(\tau\kappa)$ for every $E$, and
the remaining parameter $E$ is fixed by making the even block singular. The
operator is homogeneous of degree one in $(\alpha,\beta,\gamma)$;
its essential parameters are $\tau=\beta/\alpha\in(0,1)$
and $a=\gamma/\alpha>0$. Write
$\rho_{\tau,a}=\rho(1,\tau,a)$ and
$\rho_\tau=\rho_{\tau,1}$, with coefficient matrix $X_\tau$. These symbols denote
\emph{unnormalised} operators, not the normalised state $\rho_t$ and matrix $X_t$ of Section~\ref{sec:ppt3-family}; division by their trace gives physical density matrices. Positive scaling
does not change PPT, rank or separability.

\begin{lemma}[a $\tau$-independent supported witness]\label{lem:ppt3-h0}
\cls{proved analytically; SW-7}
Let $H_0$ be the Hermitian form on $\mathbb C^5$ in the basis $(u,v,w,e,o)$ with diagonal
$(-24-8\sqrt3,\ 1,\ 9,\ 12+4\sqrt3,\ 12+12\sqrt3)$ and $(H_0)_{vw}=(H_0)_{wv}=-\nu$, $\nu=7+4\sqrt3$. Put
$y_1=3+2\sqrt3$, $x_0=2\sqrt3-3$, $p=26+16\sqrt3$, $q=\nu^2$ (so $x_0^2q=y_1^2$), and for $z=x+iy$ let $\xi=x^2$,
$\zeta=y^2$, $r=\xi+\zeta$. Then, with $q_1$ of \eqref{eq:ppt3-curve},
\begin{equation}\label{eq:ppt3-h0}
r\,q_1\dg H_0q_1=\zeta(r-1)^2(r-y_1)^2+\xi(r-x_0)^2(r^2+pr+q)+16\nu\,\xi\zeta r\ \ge0 ,
\end{equation}
so $H_0$ is nonnegative on all cut-product rays of $R$, and $\tr(H_0X)\ge0$ whenever $BXB\dg$ is biseparable.
\end{lemma}
\emph{Idea of the proof.} At fixed $r=|z|^2$ the form $q_1\dg H_0q_1$ is a quadratic in $\xi=x^2$, so it is fixed by its two axis values
($z$ imaginary and $z$ real) and its leading coefficient. Each axis value is a perfect square times a manifestly nonnegative factor.
The computation is in Appendix~\ref{sec:ppt3-app-surface}.

\begin{theorem}[fully-PPT GME states of profile $(5,5,5,5)$]\label{thm:ppt3-orbit}
\cls{proved analytically; FP-1, SW-4, SW-5, SW-6, SW-7}
For $1/\sqrt3<\tau<1$ and every $a>0$, the normalised
$\rho_{\tau,a}$ is fully PPT and GME, with rank profile
$(5,5,5,5)$. The family has the following structure.
\begin{enumerate}[label=(\alph*)]
\item Each operator is cyclically and parity invariant,
with range contained in $R_{1/a}$, and its partial
transposes have rank at most five. In the stated interval
the range equals $R_{1/a}$.
\item For $F_\lambda=\operatorname{diag}(1,\lambda)^{\otimes3}$,
\[
\rho(\alpha,\beta,a\alpha)
=aF_\lambda\,\rho(\alpha,\beta,\alpha)F_\lambda^\dagger,
\qquad \lambda^2=a^{-1}.
\]
Thus varying $a$ is an invertible local filter;
PPT and GME depend only on $\tau$.
\item The window is sharp within this ansatz:
$\rho_{\tau,a}$ is not positive for $0<\tau<1/\sqrt3$.
At $\tau=1/\sqrt3$ it has rank four and is separable
across every cut. As $\tau\uparrow1$ its normalisation
tends to a fully separable state
(Remark~\ref{rem:ppt3-endpoint}).
\end{enumerate}
The GME proof uses the fixed supported witness $H_0$
of Lemma~\ref{lem:ppt3-h0} when $1/\sqrt3<\tau\le3/4$, and
$H_5$ from Proposition~\ref{prop:ppt3-witness}
when $3/4\le\tau<1$; the case $a\ne1$ follows because the local filter of (b)
preserves biseparability, full PPT and ranks.
\end{theorem}

\emph{Idea of the proof.} Homogeneity and the local filter reduce the problem to $\alpha=\gamma=1$. With
$Q=\tau+\sqrt{1+\tau^2}$, the odd parity block of $\rho_\tau^{T_A}$ has an explicit Gram factorisation of rank two, and the even block
has vanishing determinant and a third leading minor that changes sign at $\tau=1/\sqrt3$. Together with positivity of
$\rho_\tau$, these give the sharp window and ranks; cyclic symmetry covers the other cuts. The two witnesses cover
the window by one-variable sign checks.

\begin{corollary}[an explicit extreme state with rational weights and profile $(5,5,5,5)$]\label{cor:ppt3-extreme5}
\cls{exact-rational certificate, computer-assisted; the witness proof of genuine entanglement is exact-symbolic; extremality on the whole surface is proved from the $2\times4$ low-rank theorem and independently certified by exact-symbolic block minors; FP-2, FP-3, FP-4, FP-5, SW-3, DS-6, SB-1, SB-2}
Let $\tilde u=2\ket{000}-\ket{011}-\ket{101}-\ket{110}$, let $v,w,e,o$ be as in \eqref{eq:ppt3-vectors}, and
\begin{multline}\label{eq:ppt3-sigma5}
N=825000\,\tilde u\tilde u\dg+1440000\,ww\dg+2021184\,vv\dg\\
+145200\,(wv\dg+vw\dg)+951720\,ee\dg+450625\,oo\dg,
\qquad\sigma=\frac N{16323219}.
\end{multline}
Then $\sigma$ has rank five and is invariant under cyclic permutations and global parity; $N=30000\,\rho(55,48,110)$ is
the member $(\alpha,\beta,\gamma)=(55,48,110)$ of Theorem~\ref{thm:ppt3-orbit}, the Pythagorean choice $(m,n,a)=(8,3,2)$ with
$\alpha=m^2-n^2$, $\beta=2mn$ and $\gamma=a\alpha$.
Here $\tau=48/55$ and $s=1/2$; all six coefficients
in \eqref{eq:ppt3-rho-surface} are rational.
The projector coefficients are rational, but the matrix entries of $\sigma$ in the computational basis
belong to $\mathbb Q(\omega)$, since $e$ and $o$ contain cube-root phases.
All three partial transposes of $\sigma$ are positive semidefinite of rank five (at $A|BC$ the even and odd
blocks have ranks $3$ and $2$). The state $\sigma$ is an extreme point of the set of fully-PPT three-qubit states
and is GME, and $(1-p)\sigma+pI/8$ is GME for every $p<4.05\times10^{-3}$ (the certificate is in Appendix~\ref{sec:ppt3-app-surface}, using
the witness of Corollary~\ref{cor:ppt3-noise-reused}). More generally, \emph{every} normalised $\rho_{\tau,a}$ of Theorem~\ref{thm:ppt3-orbit}
($1/\sqrt3<\tau<1$, $a>0$) is an extreme point of the set of fully-PPT three-qubit states, and an edge state across every cut: no product
vector across a cut lies in its range with its partial conjugate in the range of the corresponding partial transpose.
\end{corollary}

\emph{Idea of the proof.} Theorem~\ref{thm:ppt3-orbit} gives positivity and ranks, and $H_{1/2}$ detects GME directly.
The noise bound uses the level-two witness $Y_*$ of Corollary~\ref{cor:ppt3-noise-reused}. For extremality, a $2\times4$ PPT state of birank $(5,5)$ that is entangled across a cut is extreme in that cut's PPT set, because a state on the boundary of its face has rank or partial-transpose rank at most four
and is then separable~\cite{KCKL,HLVC}; independently, the cyclic, parity and
reflection symmetries split the rank-$24$ face-direction calculation into blocks whose minors are explicit polynomials in
$Q=\tau+\sqrt{1+\tau^2}$ with no root in the window (Appendix~\ref{sec:ppt3-app-surface}).

\begin{proposition}[the profile-$(5,5,5,5)$ states are not isolated]\label{prop:ppt3-moduli}
\cls{exact-symbolic certificate, computer-assisted; the smoothness and openness steps are proved from the exact ranks; MO-1, MO-2, MO-3, MO-4, SB-1, SB-3, SB-4}
Near every state $\rho$ of Theorem~\ref{thm:ppt3-orbit} (in particular near $\sigma$ of Corollary~\ref{cor:ppt3-extreme5}), the fully-PPT three-qubit states of rank profile $(5,5,5,5)$ form
a smooth real manifold $\mathcal M$ of dimension $27$, and every state of $\mathcal M$ near $\rho$ is GME and an extreme
point of the set of fully-PPT states. Invertible local filters ($\rho\mapsto F\rho F\dg/\tr(F\rho F\dg)$,
$F=F_1\otimes F_2\otimes F_3$) act on $\mathcal M$ with orbits of dimension $18$, so $\mathcal M$ has nine local
moduli. The construction of Theorem~\ref{thm:ppt3-orbit} is only a two-parameter
cyclic-symmetric family (one parameter modulo filtering), not a classification.
\end{proposition}

\emph{Idea of the proof.} Rank-five states have $54$ real parameters. Each rank-five partial transpose imposes
$(8-5)^2=9$ smooth equations; the symmetry blocks of their differentials have explicit minors in $Q$ that do not vanish on the window, so they are
independent at every state of the surface. The implicit function theorem gives dimension $27$, and local filters remove
$18$ parameters (the stabiliser of the support is finite, Corollary~\ref{cor:ppt3-chirality}). GME and extremality persist nearby
(Appendix~\ref{sec:ppt3-app-surface}).

\begin{remark}[the two endpoints; two witnesses are needed]\label{rem:ppt3-endpoint}
\cls{proved; endpoints and uniqueness checked symbolically in SW-7, SW-8}
At $\tau=1/\sqrt3$ ($Q=\sqrt3$) the operator $\rho_\tau$ is a limit of fully-PPT states, hence positive semidefinite with all three
partial transposes positive semidefinite, and it has rank four, so it is separable across each cut by the $2\times4$
low-rank theorem~\cite{KCKL,HLVC}. Explicitly, at $a=1$ its coefficient matrix
$X_{1/\sqrt3}=\operatorname{diag}(1,3\sqrt3,\sqrt3/3,2,2\sqrt3/3)+\sqrt3\,(E_{vw}+E_{wv})$ is a positive combination of four
$A|BC$-product rays, at the zeros $z=\pm i\sqrt{y_1},\pm\sqrt{x_0}$ of $q_1(z)\dg H_0q_1(z)$ in~\eqref{eq:ppt3-h0}:
\begin{equation}\label{eq:ppt3-lower-endpoint}
 X_{1/\sqrt3}
  =\frac{3\sqrt3-5}{48}\sum_{\pm}q_1\bigl(\pm i\sqrt{y_1}\bigr)q_1\bigl(\pm i\sqrt{y_1}\bigr)\dg
  +\frac{3\sqrt3+5}{48}\sum_{\pm}q_1\bigl(\pm\sqrt{x_0}\bigr)q_1\bigl(\pm\sqrt{x_0}\bigr)\dg ,
\end{equation}
with weights $(3\sqrt3\mp5)/48>0$. Cyclic permutation gives the decompositions across $B|CA$ and $C|AB$, and the local filter in
Theorem~\ref{thm:ppt3-orbit}(b) carries them to every $a>0$.
It is not fully separable. Writing the product vectors of Proposition~\ref{prop:ppt3-witness}(a) as
$B_sq_s(z)=(\ket0+z\ket1)\otimes\phi_s(z)$, the $BC$ coefficient matrix of $\phi_s(z)$ has determinant
$-3(z^2+s)^2$, so the only fully product rays in $R_s$ are $\ket{111}$ and $(\ket0\pm i\sqrt s\ket1)^{\otimes3}$.
A fully separable state with range in $R_{1/a}$ therefore has rank at most three, while $\rho_{1/\sqrt3,a}$ has rank four.

At the upper endpoint, for every fixed $a>0$ the limit is fully separable and has rank profile $(2,2,2,2)$.
Indeed the weights satisfy $(O,E,G,V)\to(0,0,-1/a,1/a^2)$ as $\tau\uparrow1$, giving
\[
 \rho_{1,a}:=\lim_{\tau\uparrow1}\rho_{\tau,a}
    =a\,u_{1/a}u_{1/a}\dg+(w-a^{-1}v)(w-a^{-1}v)\dg.
\]
Put $\ket{\eta_\pm(a)}=(\sqrt a\ket0\pm i\ket1)/\sqrt{a+1}$. Expanding the product vectors gives
\begin{equation}\label{eq:ppt3-upper-endpoint}
 \frac{\rho_{1,a}}{\tr \rho_{1,a}}
   =\frac12\left[
     (\ket{\eta_+(a)}\bra{\eta_+(a)})^{\otimes3}
    +(\ket{\eta_-(a)}\bra{\eta_-(a)})^{\otimes3}\right],
 \qquad \tr \rho_{1,a}=\frac{(a+1)^3}{a^2}.
\end{equation}
The two local rays are distinct for every finite $a>0$, and remain distinct after conjugating any subset of
factors. Thus both the state and its three partial transposes have rank exactly two. The GME interval therefore
terminates at states separable across every cut at both ends; only the upper endpoint is fully separable.

Together, \eqref{eq:ppt3-lower-endpoint} and~\eqref{eq:ppt3-upper-endpoint} explain the origin of $H_0$ and show that no
single witness detects $\rho_\tau$ on the whole window $1/\sqrt3<\tau<1$, so two witnesses are needed. Only the compression of a
witness to $R$ matters, and it is a form $H$ on $\mathbb C^5$ nonnegative on the cut-product rays of $R$. If $\tr(HX_\tau)<0$
on the whole window, continuity gives $\tr(HX_{1/\sqrt3})\le0$ and $\tr(HX_1)\le0$, where $X_1$ is the coefficient matrix of
$\rho_{1,1}$. Both endpoint matrices are positive
combinations of cut-product rays, so $q_1(z)\dg Hq_1(z)$ vanishes, and being nonnegative has a double zero, at
$z=\pm i\sqrt{y_1},\pm\sqrt{x_0}$ and at $z=\pm i$. The cyclic, parity and $K\circ S_{BC}$ symmetries fix every $X_\tau$ and
permute the cut-product rays, so the average $\bar H$ of $H$ over them has the same properties. These double zeros determine
$\bar H$ up to a nonnegative multiple of $H_0$ (Appendix~\ref{sec:ppt3-app-surface}). But $\tr(H_0X_\tau)$ changes sign at
$\tau\approx0.805$, the root of $P(Q)$ (for instance $\tr(H_0X_{48/55})>0$), a contradiction.
\end{remark}

\begin{remark}[the edge condition fails exactly at the endpoint rays]\label{rem:ppt3-edge-overlap}
\cls{exact-symbolic certificate, computer-assisted; EO-1, EO-2, EO-3, EO-4, EO-5}
Take $a=1$; invertible local filters map product vectors to product vectors and commute with partial conjugation up to the conjugate
filter, so they preserve the edge property. By Proposition~\ref{prop:ppt3-witness}(a) the $A|BC$ product rays of $R_1$ are those of
$\psi(z)=Bq_1(z)=(\ket0+z\ket1)\otimes\phi_1(z)$, $z\in\mathbb C$, and of $v=\ket{111}$. So $\rho_\tau$ is an edge state across $A|BC$ exactly when none of them has its partial conjugate
$\psi^{T_A}(z)=(\ket0+\bar z\ket1)\otimes\phi_1(z)$ in $\operatorname{ran}\rho_\tau^{T_A}$, that is, orthogonal to the kernel vectors $k_1,k_2,k_3$ of Appendix~\ref{sec:ppt3-app-surface}.
For $z=x+iy$,
\[
k_1\dg\psi^{T_A}(z)=6\bigl((Q^2-1)|z|^2+Q(z^2-1)\bigr),\qquad\operatorname{Im}=12Qxy .
\]
A zero therefore has $y=0$, $x^2=Q/(Q^2+Q-1)$, or $x=0$, $y^2=Q/(Q^2-Q-1)$, and there
\[
k_3\dg\psi^{T_A}=-\frac{6Qx(Q^2-3)(Q^2+2Q-1)}{Q^2+Q-1}\quad\text{resp.}\quad-\frac{6iQy(Q^2-3)(Q^2-2Q-1)}{Q^2-Q-1},
\]
while $k_2\dg\ket{111}=12$. On $\sqrt3<Q<1+\sqrt2$ no factor vanishes, so the edge condition holds; this checks the edge property across $A|BC$ directly, and across the other cuts by
the cyclic shift. At the endpoints the factors $Q^2-3$ and $Q^2-2Q-1$ vanish. At $Q=\sqrt3$ the two branches give $z=\pm\sqrt{x_0},\pm i\sqrt{y_1}$, the four rays of~\eqref{eq:ppt3-lower-endpoint}, and at $Q=1+\sqrt2$
the imaginary branch gives $z=\pm i$, the rays of~\eqref{eq:ppt3-upper-endpoint}. So at both ends the common zeros of the overlaps with $k_1,k_2,k_3$ are exactly the product rays of the separable decomposition of the limit state.
\end{remark}


\section{The states are in the mixed \texorpdfstring{$W$}{W} class}\label{sec:ppt3-w-class}

Pure three-qubit GME states fall into two inequivalent classes under stochastic local operations, represented by the GHZ
and $W$ states~\cite{DVC}; the GHZ class is distinguished by a nonzero three-tangle~\cite{CKW}, the modulus of Cayley's
hyperdeterminant $\operatorname{Det}$ (up to a factor four). Ac\'in et al.\ extended this to mixed states and conjectured that PPT entangled three-qubit states
lie in the mixed $W$ class~\cite{ABLS}. Our states contain the GHZ-type vector $u$ in their spectral decomposition, but
this does not make them GHZ-class: other decompositions can avoid GHZ-type components altogether, and we construct them.
Write $\mathcal W$ for the convex hull of pure three-qubit states with vanishing Cayley hyperdeterminant,
including biseparable states. This is the mixed $W$ class of~\cite{ABLS}; its pure generators have zero three-tangle.
An invertible local filter preserves this class, since the hyperdeterminant transforms by the squared product
of the local determinants. Convex mixing and local quantum channels also preserve $\mathcal W$.

\begin{theorem}[the constructed fully-PPT states belong to the mixed $W$ class]\label{thm:ppt3-w-class}
\cls{proved analytically; WC-1, WC-2, WC-3, WC-4}
Every normalised state $\rho_t$ of Theorem~\ref{thm:ppt3-family} in its PPT interval, the extreme state
$\sigma_6$ of Corollary~\ref{cor:ppt3-extreme6}, every normalised $\rho_{\tau,a}$ of
Theorem~\ref{thm:ppt3-orbit} ($1/\sqrt3<\tau<1$, $a>0$), the phase-coexistence line $\sigma(t)$ of
Proposition~\ref{prop:ppt3-exposed-transition} ($-1/200\le t\le1$), and the full-rank state $\tilde\rho$ of
Corollary~\ref{cor:ppt3-robust} belong to $\mathcal W$. So do their local-filter images and their mixtures
with fully separable noise. In their GME ranges these states lie in $\mathcal W\setminus\mathrm{BS}$.
\end{theorem}

\begin{proof}
All vectors below are unnormalised. Put $[x]=xx\dg$ and let $\mathcal C$ average the three cyclic permutations,
while $\mathcal P(X)=(X+Z^{\otimes3}XZ^{\otimes3})/2$ averages parity. These maps preserve $\mathcal W$.
On either parity subspace only four computational basis states occur, and Cayley's hyperdeterminant reduces to
\[
 \operatorname{Det}\Bigl(\sum_{\text{even}}c_{ijk}\ket{ijk}\Bigr)=4c_{000}c_{011}c_{101}c_{110},\qquad
 \operatorname{Det}\Bigl(\sum_{\text{odd}}c_{ijk}\ket{ijk}\Bigr)=4c_{111}c_{100}c_{010}c_{001}.
\]
So a vector of definite parity with one vanishing amplitude has zero hyperdeterminant. This applies to
$v,w,e,\bar e,o,\bar o$; to the even vectors $u_s+xe+y\bar e$, where $u_s=\ket{000}-s(\ket{011}+\ket{101}+\ket{110})$ and
$x,y$ are real with $x+y=s$ (the amplitude of $\ket{011}$ is $-s+x+y=0$); and to the odd vectors $v+b(w-o)$, $b\in\mathbb C$
(the amplitude of $\ket{100}$ vanishes). For vectors mixing the two parities, direct substitution gives
\[
 \operatorname{Det}(u+ze)=4(z^3-1),\qquad
 \operatorname{Det}(u+zv)=z^2-4,\qquad
 \operatorname{Det}\bigl(xe+i(v-\beta w+\gamma o)\bigr)=4(\gamma^3-\beta^3+3\beta\gamma x^2),
\]
so for real $0<\gamma\le\beta$ the last vector has zero hyperdeterminant when $x^2=(\beta^3-\gamma^3)/(3\beta\gamma)$.
Also $u+i(w-v)=(\ket0+i\ket1)^{\otimes3}$ is a product vector. Hence the following operators lie in the cone over $\mathcal W$:
\begin{equation}\label{eq:ppt3-w-atoms}
\begin{aligned}
 \mathcal C([u_s+xe+y\bar e])&=[u_s]+x^2[e]+y^2[\bar e]\quad(x,y\in\mathbb R,\ x+y=s),\\
 \mathcal C([v+b(w-o)])&=[v+bw]+|b|^2[o]\quad(b\in\mathbb C),\\
 \mathcal P([u+2v])&=[u]+4[v],\qquad \mathcal P([u+i(w-v)])=[u]+[w-v],\\
 \mathcal C\bigl(\mathcal P([xe+i(v-\beta w+\gamma o)])\bigr)&=[v-\beta w]+\gamma^2[o]+\frac{\beta^3-\gamma^3}{3\beta\gamma}[e]
 \quad(0<\gamma\le\beta).
\end{aligned}
\end{equation}
For $s=1$, $x=1$, $y=0$ the first is $\mathcal C([u+e])=[u]+[e]$.

First consider a positive diagonal coefficient matrix
$d_u[u]+d_v[v]+d_w[w]+d_e[e]+d_o[o]$, with all five weights nonnegative.
If $d_u\le d_e+d_v/4$, put $r=\min(d_u,d_e)$ and $q=d_u-r$. It equals
\[
 r\,\mathcal C([u+e])+q\,\mathcal P([u+2v])
 +(d_e-r)[e]+(d_v-4q)[v]+d_w[w]+d_o[o],
\]
with nonnegative coefficients. The first family satisfies $d_u=t$, $d_v=3t$, $d_e=3$ and $t\le t_+<4$,
so $t\le3+3t/4$. For $\sigma_6$, $d_u=t_-<4$, $d_e=3$ and $d_v=8(9-t_-)/5>8$, which also suffice.

For the surface, the filter law in Theorem~\ref{thm:ppt3-orbit}(b) reduces the problem to $a=1$.
Use the notation of Appendix~\ref{sec:ppt3-app-surface}: the weights of $u$ and $w$ are $1$ and $\tau$, and
\[
 \kappa=1-\tau^2,\quad Q=\tau+\sqrt{1+\tau^2},\quad
 E=\kappa Q^2,\quad O=\kappa/\tau,\quad G=(E-1)/\tau.
\]
and $V=(E^2-2\kappa E+\kappa)/(\tau\kappa)$. The coefficient block on $(v,w)$ is positive semidefinite, since the operator
is, so $V>0$ and $\tau V-G^2\ge0$.
Its other useful identity is
\begin{equation}\label{eq:ppt3-w-budget}
 OV-G^2=2E-1.
\end{equation}
If $E\ge1$, put $b=G/V$. Then the surface operator is
\[
 \mathcal C([u+e])+V\mathcal C([v+b(w-o)])
 +(\tau-G^2/V)[w]+(E-1)[e]+(O-G^2/V)[o].
\]
All coefficients are nonnegative, the last by~\eqref{eq:ppt3-w-budget}.

If $E<1$, put $\mu=(1-E)/\tau=-G$. Since $Q\ge1$ and $0<\tau<1$, $E=\kappa Q^2\ge\kappa=1-\tau^2\ge1-\tau$, so $0<\mu\le1$ and
\[
 \tau-\mu=\frac{E-\kappa}{\tau}\ge0,\quad
 V-\mu=\frac{E(E-\kappa)}{\tau\kappa}\ge0,\quad
 E-1+\mu=(1-E)(\tau^{-1}-1)\ge0.
\]
Now the surface operator is
\[
 \mu\,\mathcal P([u+i(w-v)])+(1-\mu)\mathcal C([u+e])
 +(\tau-\mu)[w]+(V-\mu)[v]+(E-1+\mu)[e]+O[o].
\]
This again is a positive sum of zero-three-tangle projectors.

On the phase-coexistence line $\sigma(t)=(1-t)\sigma_\partial+t\sigma_{\rm BS}$ of
Proposition~\ref{prop:ppt3-exposed-transition}, $\sigma(t)$ is biseparable, hence in $\mathcal W$, for $0\le t\le1$. By
convexity it remains to treat $t=-1/200$. There the coefficient matrix $(1-t)X_\partial+tX_{\rm BS}$ of
\eqref{eq:ppt3-xdelta} and~\eqref{eq:ppt3-bs-same-support} has diagonal $(d_u,d_v,d_w,d_e,d_o)$ and $(v,w)$ entry $-\mu$ with
$\mu=10930399/348000000>0$. Subtracting $\mu\,\mathcal P([u+i(w-v)])$ leaves the positive diagonal matrix with weights
$(d_u-\mu,d_v-\mu,d_w-\mu,d_e,d_o)$, and $d_u-\mu-d_e=7/31250\le(d_v-\mu)/4$, so the diagonal case above applies.

Finally, write the full-rank state $\tilde\rho$ of Corollary~\ref{cor:ppt3-robust}, displayed in~\eqref{eq:ppt3-tilde-rho},
as
\[
 10^7\tilde\rho=c_0[g]+c_1(gh\dg+hg\dg)+c_2[h]+c_e[e]+c_{\bar e}[\bar e]+c_v[v]+c_{vw}(vw\dg+wv\dg)+c_w[w]+c_o[o]+c_{\bar o}[\bar o],
\]
with $g=\ket{000}$, $h=\ket{011}+\ket{101}+\ket{110}$ and integer coefficients. The $2\times2$ blocks on $(g,h)$ and
$(v,w)$ have positive determinants $c_0c_2-c_1^2=382\,804\,344$ and $c_vc_w-c_{vw}^2=257\,058\,514$. With
$s=-c_1/c_0=232675/238898\in(0,1)$ and $\beta=-c_{vw}/c_v=520181/390647\in(1,4/3)$, completing the square writes them as
$c_0[u_s]+(c_2-c_0s^2)[h]$ and $c_v[v-\beta w]+(c_w-c_v\beta^2)[w]$, both remainders positive. Take $(x,y)=(\tfrac59s,\tfrac49s)$
in the first atom of~\eqref{eq:ppt3-w-atoms}, and $\gamma=1$, $x'^2=(\beta^3-1)/(3\beta)<37/108$ in the last. Then
\begin{multline*}
 10^7\tilde\rho=c_0\,\mathcal C([u_s+xe+y\bar e])+c_v\,\mathcal C\bigl(\mathcal P([x'e+i(v-\beta w+o)])\bigr)\\
 +(c_2-c_0s^2)[h]+r_e[e]+r_{\bar e}[\bar e]+(c_w-c_v\beta^2)[w]+(c_o-c_v)[o]+c_{\bar o}[\bar o],
\end{multline*}
where $r_{\bar e}=c_{\bar e}-c_0y^2>c_{\bar e}-\tfrac{16}{81}c_0>0$, $r_e=c_e-c_0x^2-c_vx'^2>c_e-\tfrac{25}{81}c_0-\tfrac{37}{108}c_v>0$
and $c_o-c_v=165\,771>0$.
Dividing by the trace converts every displayed conic decomposition to a convex decomposition of normalised states.
\end{proof}

\begin{remark}[scope of the classification]\label{rem:ppt3-w-class-scope}
This settles the mixed-$W$-class question for every explicit state constructed in this paper. It also gives fully-PPT GME
states in $\mathcal W$ with nonempty interior: states separable across $A|BC$ lie in $\mathcal W$ and contain a ball around
$I/8$~\cite{GB}, so for $0<p\le2\times10^{-5}$ the mixtures $(1-p)\tilde\rho+p\sigma$ with $\sigma$ in that ball form a ball
in $\mathcal W$, all within trace distance $p$ of $\tilde\rho$ and hence fully PPT and GME by Corollary~\ref{cor:ppt3-robust}.
It does not settle the conjecture of~\cite{ABLS} for arbitrary fully-PPT three-qubit states. Nor does a zero-three-tangle
decomposition imply biseparability: these states are GME in their stated ranges.
\end{remark}

\section{One support, both phases}\label{sec:ppt3-coexistence}

The product curves constrain biseparable decompositions, but the support alone does not decide GME. The support
$R=R_1$ of the first family carries both biseparable and GME fully-PPT states with the same rank profile.
An explicit line crosses their boundary while remaining strictly PPT. The proofs, with the exact matrices and minors, are in
Appendix~\ref{sec:ppt3-app-coexistence}.

Biseparable states on $R$ can be written down explicitly from the product curve of Proposition~\ref{prop:ppt3-witness}. For
finite $z$ let $p_z=Bq_1(z)/\|Bq_1(z)\|$, a normalised $A|BC$-product vector in $R$, let $\Pi$ be the cyclic qubit permutation,
and define the cyclically and parity symmetrised mixture
\begin{equation}\label{eq:ppt3-gamma}
 \Gamma_z=\frac16\sum_{j=0}^2\sum_{\epsilon=\pm1}
   \Pi^j\ket{p_{\epsilon z}}\bra{p_{\epsilon z}}
   \Pi^{-j}.
\end{equation}
Since $Z^{\otimes3}\ket{p_z}=-\ket{p_{-z}}$, the sum over $\epsilon=\pm1$ is the global parity average, and each $\Gamma_z$
is an equal mixture of six pure product states across the three cuts, so $\Gamma_z$ and its convex mixtures are biseparable and supported on $R$.

\begin{theorem}[biseparable and GME states on one PPT support]
\label{thm:ppt3-phase-coexistence}
\cls{proved; exact rational arithmetic, PG-2}
The fixed support $R=R_1$ contains both a
biseparable state and a GME state of rank five,
\emph{each} PPT with strictly positive partial
transposes. They therefore have the same rank
profile $(5,8,8,8)$ and exactly the same
three cut-product curves. An explicit
biseparable example is
$\sigma_{\rm BS}=\frac13\bigl(\ket{111}\bra{111}+\Gamma_2+\Gamma_{2i}\bigr)=BX_{\rm BS}B^\dagger$,
where the rational coefficient matrix $X_{\rm BS}$ in the basis $(u,v,w,e,o)$ is given in
\eqref{eq:ppt3-bs-same-support} of Appendix~\ref{sec:ppt3-app-coexistence}.
The GME example is $\rho_{10/3}$ from
Theorem~\ref{thm:ppt3-family}.
Within the trace-one operators supported on $R$,
the fully-PPT biseparable states and the fully-PPT
GME states both have nonempty relative interiors.
\end{theorem}

\emph{Idea of the proof.} The state $\sigma_{\rm BS}$ is biseparable by construction; expanding the mixture gives
$X_{\rm BS}$, and Sylvester's criterion on the two parity blocks of one partial transpose gives strict PPT. The state
$\rho_{10/3}$ has the same range and is GME by \eqref{eq:ppt3-h5-value}. Strict PPT, rank five and GME persist under small
perturbations within $R$. On the biseparable side, averaging curve projectors over a measure of full support gives an
interior point of the supported biseparable cone: no nonzero supporting form can vanish on that average. Mixing it into
$\sigma_{\rm BS}$ retains strict PPT (Appendix~\ref{sec:ppt3-app-coexistence}).

\begin{proposition}[an exposed phase boundary inside the strict-PPT region]
\label{prop:ppt3-exposed-transition}
\cls{proved; exact rational arithmetic, PG-4}
With the biseparable cut-product mixtures $\Gamma_z$ of
\eqref{eq:ppt3-gamma}, define
\begin{equation}\label{eq:ppt3-exposed-state}
 \sigma_\partial=\frac15 \Gamma_1+\frac12 \Gamma_i+
                  \frac3{10}\Gamma_{i\sqrt5}
 =BX_\partial B^\dagger ,
\end{equation}
with the rational matrix $X_\partial$ of \eqref{eq:ppt3-xdelta} in Appendix~\ref{sec:ppt3-app-coexistence}.
This rank-five biseparable state has
$\sigma_\partial^{T_c}\succ I/300$ at each cut, yet
$\tr(H_5X_\partial)=0$. The straight line
\begin{equation}\label{eq:ppt3-exact-phase-line}
 \sigma(t)=(1-t)\sigma_\partial+t\sigma_{\rm BS},
          \qquad -\frac1{200}\le t\le1,
\end{equation}
consists entirely of rank-five fully-PPT states
with strictly positive partial transposes. Along
this \emph{single fixed-support, fixed-rank} line,
\[
 \sigma(t)\text{ is biseparable for }0\le t\le1,
 \qquad
 \sigma(t)\text{ is GME for }-1/200\le t<0.
\]
Thus $t=0$ is an exact biseparable--GME transition
\emph{inside} the strict-PPT region, exposed by $H_5$.
\end{proposition}

\emph{Idea of the proof.} The factorisation \eqref{eq:ppt3-h5} vanishes at $|z|=1$ and at $z=\pm i\sqrt5$, so $H_5$ has zero
expectation on each summand of $\sigma_\partial$. Since $\tr(H_5X_{\rm BS})>0$, the witness value is linear in $t$ along
the line and negative exactly for $t<0$, which excludes biseparability there. The PPT margin is a Sylvester check, a
Hilbert--Schmidt estimate keeps it for the small negative range of $t$, and convexity covers $0\le t\le1$. The exact data are
in Appendix~\ref{sec:ppt3-app-coexistence}.

\section{Noise robustness}\label{sec:ppt3-noise}

How much white noise or independent qubit depolarisation can the states tolerate while remaining GME?
We bound the tolerance from below with exact symmetric-extension witnesses, and from above with explicit biseparable
decompositions completed by the Gurvits--Barnum separable ball~\cite{GB}. The bounds meet to within a few per cent in
relative terms. Appendix~\ref{sec:ppt3-app-noise} gives the certificate checks.

The supported witnesses establish GME on a five-dimensional support.
Mixing in a small amount of white noise yields full-rank, strictly PPT
states while preserving GME, giving the following general consequence.

\begin{corollary}[positive volume in any finite dimension]\label{cor:ppt3-volume}
\cls{proved; analytic consequence of Theorem~\ref{thm:ppt3-family}}
In any finite multipartite system, if a fully-PPT GME state
exists, such states have non-empty interior (and hence positive
volume) among density matrices. In particular this holds for
three qubits. The conclusion is analytic; the certificates below
provide explicit noise and preparation radii.
\end{corollary}
\begin{proof}
Let $\rho$ be any fully-PPT GME state on a Hilbert space of total
dimension $d$. The biseparable set $\mathrm{BS}$ is compact, being the
convex hull of the union of finitely many compact sets of pure
cut-product states. Since $\rho\notin\mathrm{BS}$, the noisy state
$\rho_p=(1-p)\rho+pI/d$ remains outside $\mathrm{BS}$ for all
sufficiently small $p>0$. But $\rho_p\succeq pI/d$ and
$\rho_p^{T_S}=(1-p)\rho^{T_S}+pI/d\succeq pI/d$ at every cut $S$.
Thus $\rho_p$ is full rank and strictly PPT, and it is GME.
Continuity gives an open neighbourhood of $\rho_p$ in the trace-one
Hermitian affine space consisting entirely of physical fully-PPT GME
states. Applying this observation to Theorem~\ref{thm:ppt3-family}
proves the three-qubit claim.
\end{proof}

Positive volume does not say how large the margins are. Since PPT tests cannot detect these states, certifying an
explicit margin requires a stronger relaxation of biseparability. We use dual witnesses
from level two of the Doherty--Parrilo--Spedalieri (DPS) symmetric-extension hierarchy~\cite{DPS02,DPS,DPS05}, the lowest level
stronger than PPT. Applying it across each cut to relax biseparability was proposed in Appendix~B of~\cite{JMG}.
Every cut-product vector extends by duplicating its qubit factor; the following certificate exploits this necessary
condition. Numerical searches supplied candidates, rounded to rational parameters; only exact checks establish the
bounds (Appendix~\ref{sec:ppt3-verifiers}).
Here $\mathrm{Sym}^2(A)$ is the subspace invariant under exchanging the two copies of $A$; tensor subscripts specify
which factors an operator acts on.

\begin{lemma}[certificate and stability principle]\label{lem:ppt3-dps}
\cls{proved}
Let $A|B$ be a bipartition into finite-dimensional factors, let $V$ have full column rank with image
$\mathrm{Sym}^2(A)\otimes B$ inside $A_1\otimes A_2\otimes B$ (the columns need not be normalised), and let $Y$ be
Hermitian on $A\otimes B$ and $Z_1,Z_2\succeq0$ on $A_1\otimes A_2\otimes B$ with
\begin{equation}\label{eq:ppt3-cert}
P=V\dg\bigl(Y_{A_1B}\otimes I_{A_2}-Z_1^{T_B}-Z_2^{T_{A_2}}\bigr)V\succeq0 .
\end{equation}
Then $\langle x|Y|x\rangle\ge0$ for every product vector $x=a\otimes b$. Consequently, on three qubits, if $Y$ is invariant
under cyclic permutations of the qubits and satisfies \eqref{eq:ppt3-cert} for the cut $A|BC$, then
$\tr(Y\sigma)\ge0$ for every biseparable state $\sigma$. If moreover $\rho$ is a three-qubit state with
$\tr(Y\rho)=-g<0$, $h=\tr(Y)/8>0$, and $\operatorname{spec}(Y)\subset[l,u]$, then
\begin{equation}\label{eq:ppt3-stability}
(1-p)\rho+pI/8\ \text{is GME for }p<\frac g{g+h},\qquad D(\rho,\mathrm{BS})\ge\frac g{u-l}.
\end{equation}
Finally, if $\rho^{T_j}\succeq cI$ at every cut $j$, with positive $c$, then every
physical state $\sigma$ with
\[
D(\rho,\sigma)<\frac{2c}{1+\sqrt2}
\]
is strictly PPT across every cut. This constant is dimension independent.
\end{lemma}

\emph{Idea of the proof.} For a product vector $x=a\otimes b$, the certificate writes $\|a\|^2\langle x|Y|x\rangle$ as the sum of a positive
semidefinite quadratic form (from $P$) and two nonnegative terms (from $Z_1$ and $Z_2$), so $Y$ is nonnegative on every cut-product vector and hence on $\mathrm{BS}$.
The noise and trace-distance bounds follow from the sign and spectral width of $Y$. The PPT-stability radius uses that the partial transpose of
a pure state has spectral width at most $(1+\sqrt2)/2$. Details in Appendix~\ref{sec:ppt3-app-noise}.

Local depolarisation means applying $\mathcal D_{q_j}(X)=(1-q_j)X+q_j\tr(X)I/2$ on qubit $j$, with
\emph{independent rates} $0\le q_j\le q$. White noise means $(1-p)\rho+pI/8$.

\begin{corollary}[noise and preparation radius of the family]\label{cor:ppt3-noise}
\cls{exact-rational certificate, with proved PPT stability; DS-1, DS-2, DS-3, DS-6}
The rational, symmetry-reduced cyclic witnesses $Y_t$ of the ancillary certificate data give the following conservative tolerances, where the balls are trace-distance balls
around $\rho_t$.
\begin{center}
\small\setlength{\tabcolsep}{5pt}
\begin{tabular}{@{}lcccc@{}}\toprule
$t$ & white noise $p$ & local $q$ per qubit &
 GME ball & fully-PPT ball\\\midrule
$327/100$ & $3.3\times10^{-4}$ & $1.4\times10^{-4}$ &
 $1.0\times10^{-4}$ & $1.01\times10^{-4}$\\
$10/3$ & $8.7\times10^{-4}$ & $3.7\times10^{-4}$ &
 $2.7\times10^{-4}$ & $5.7\times10^{-4}$\\
$169/50$ & $1.4\times10^{-3}$ & $6.1\times10^{-4}$ &
 $4.4\times10^{-4}$ & $2.0\times10^{-4}$\\\bottomrule
\end{tabular}
\end{center}
Specifically, white noise
$(1-p)\rho_t+pI/8$ remains GME for $p$ up to the listed
value, and independent local depolarisation with rates
$q_j\le q$ does likewise. Every physical state in a listed closed GME ball is GME, and every
state strictly inside a listed PPT ball is fully PPT. Every state at trace distance
strictly below the smaller of the last two entries therefore has \emph{both} properties:
below $2.7\times10^{-4}$ for $t=10/3$, below $1.0\times10^{-4}$ for $t=327/100$ (where the
whole GME ball is also PPT) and below $2.0\times10^{-4}$ for $t=169/50$.
In particular $D(\rho_{10/3},\mathrm{BS})>2.7\times10^{-4}$.
White noise itself preserves PPT automatically.

One fixed dual witness also certifies \emph{uniformly}
over the exact interval \eqref{eq:ppt3-interval}:
$p\le2.4\times10^{-4}$ of white noise and a GME ball
of radius $7.5\times10^{-5}$. This gives a second GME proof, independent of the analytic witness.
\end{corollary}

\emph{Idea of the proof.} Exact witnesses satisfy Lemma~\ref{lem:ppt3-dps}. Unequal depolarisation rates reduce to one
uniform rate because further local noise preserves biseparability. Affine dependence on $t$ reduces the uniform
interval bound to two endpoint checks (Appendix~\ref{sec:ppt3-app-noise}).

\begin{corollary}[noise bounds for the surface state from one exact witness]
\label{cor:ppt3-noise-reused}
\cls{exact-symbolic certificate; DS-6}
Let $\rho_*$ be the trace-one normalisation of the analytic
surface operator $\rho(312,266,603)$ in
\eqref{eq:ppt3-rho-surface}; here
$\tau=133/156$ and $a=201/104$.
It is fully PPT and GME, with rank profile $(5,5,5,5)$.
A single exact level-two witness $Y_*$ (optimised numerically for $\rho_*$ under a spectral bound)
certifies:
\begin{enumerate}[label=(\roman*)]
\item The states $(1-p)\rho_*+pI/8$ remain fully PPT
and GME for $0\le p\le2153/500000$
(about $0.43\%$ white noise).
\item Independent local depolarisation preserves both
properties whenever each rate is at most
$81/50000$ (about $0.16\%$ per qubit).
\item Let
\[
 \rho_\circ=\left(1-\frac{32}{10000}\right)\rho_*
                  +\frac{32}{10000}\frac I8 .
\]
Every physical state $\sigma$ with
$D(\sigma,\rho_\circ)<331/1000000$
is fully PPT and GME. Thus a full-rank state
has a simultaneous preparation radius of
$3.31\times10^{-4}$.
\item Every physical state $\sigma$ with
$D(\sigma,\rho_*)\le129/100000$
is GME (a GME-only radius of $1.29\times10^{-3}$).
\end{enumerate}
\end{corollary}

\emph{Idea of the proof.} The certified expectation and spectral bounds of $Y_*$ give each margin by
Lemma~\ref{lem:ppt3-dps}; white noise opens a positive PPT gap for (iii) (Appendix~\ref{sec:ppt3-app-noise}).

\begin{corollary}[a full-rank fully-PPT GME state tolerating $2.3\%$ white noise]\label{cor:ppt3-robust}
\cls{exact-rational certificate, with proved PPT stability; DS-4, DS-6}
The three-qubit density matrix $\tilde\rho$ with seven-decimal rational coordinates in the certificate data (entries in $\mathbb Q(\sqrt{-3})$) has
full rank, is invariant under cyclic permutations of the qubits and under global parity $Z^{\otimes3}$, has every partial
transpose strictly above $2.49\times10^{-5}$ times the identity, and is GME: its rational level-two witness is
nonnegative on every cut-product vector and negative on $\tilde\rho$. Certified tolerances: global white noise
$p\le2.33\times10^{-2}$, independent depolarising noise with every rate at most $1.0\times10^{-2}$, and every state
within trace distance $7.3\times10^{-3}$ of $\tilde\rho$ is GME. The ball of radius $2.0\times10^{-5}$ is also certified
fully PPT, by the verified lower bound $2.49\times10^{-5}$ on the partial transposes and Lemma~\ref{lem:ppt3-dps}.
These are certified lower bounds; Proposition~\ref{prop:ppt3-noise-sharp} gives the matching upper bound for white noise.
\end{corollary}

\emph{Idea of the proof.} Lemma~\ref{lem:ppt3-dps} applied to the rational data of $\tilde\rho$, with exact checks of full rank, symmetry and the
partial-transpose eigenvalue bounds. Details in Appendix~\ref{sec:ppt3-app-noise}.

The state $\tilde\rho$ is obtained from the white-noise mixture $(1-\epsilon)\rho_{\rm opt}+\epsilon I/8$, $\epsilon=2.5\times10^{-4}$, of a
numerically optimised fully-PPT state $\rho_{\rm opt}$ (Remark~\ref{rem:ppt3-robust-structure}) by rounding to rational
entries; white noise provides strict PPT before rounding. Its GME ball is more than five times the GME-only radius of
Corollary~\ref{cor:ppt3-noise-reused}(iv). Its exact coordinates are displayed in Appendix~\ref{sec:ppt3-app-noise}.

For arbitrary physical perturbations, the certified GME radius is $7.3\times10^{-3}$ in trace distance, whereas the
simultaneous PPT radius is only $2.0\times10^{-5}$ because the smallest partial-transpose eigenvalue is small.
The much larger white-noise tolerance therefore does not guarantee stability under generic preparation errors.

\begin{proposition}[the white-noise thresholds are essentially sharp]\label{prop:ppt3-noise-sharp}
\cls{exact-rational/algebraic certificate, computer-assisted; DS-7}
Let $p_{\rm GME}(\rho)=\sup\{p\in[0,1]:(1-p)\rho+pI/8\text{ is GME}\}$, with $p_{\rm GME}(\rho)=0$ for biseparable $\rho$. Then
\begin{align*}
 \frac{2153}{500000}&\le p_{\rm GME}(\rho_*)\le \frac{4319}{1000000},\\
 \frac{87}{100000}&\le p_{\rm GME}(\rho_{10/3})\le \frac{89}{100000},\\
 \frac{233}{10000}&\le p_{\rm GME}(\tilde\rho)\le \frac{5839}{250000},
\end{align*}
with $\rho_*$ and $\tilde\rho$ as in Corollaries~\ref{cor:ppt3-noise-reused} and~\ref{cor:ppt3-robust}. The noisy state is
biseparable at, and for every $p$ above, each upper bound.
\end{proposition}

\emph{Idea of the proof.} The lower bounds are the three corollaries above. For the upper bounds, an explicit separable-across-$A|BC$ operator,
assembled from three orbit representatives and symmetrised cyclically, is subtracted from the noisy state; the residual lies in the
Gurvits--Barnum ball around the identity~\cite{GB}, hence is separable, so the noisy state is biseparable. Details in Appendix~\ref{sec:ppt3-app-noise}.

The certified intervals have relative widths of at most $2.3\%$, so the small margins reflect the states, not the
relaxation. Numerically the agreement is closer: for all three states the optimal level-two bound and the numerical
value of the optimised three-orbit decomposition agree to within $10^{-7}$ (numerical evidence only); the decompositions use three orbits of qubit directions each.

To guide the search for robust states, we use the convex white-noise robustness $\operatorname{rob}$ below,
rather than $p_{\rm GME}$ itself.

\begin{proposition}[an extremal reduction for white-noise robustness]\label{prop:ppt3-noise-extremal}
\cls{proved; general convex duality}
In any finite multipartite system of dimension $d$, let $\mathcal P$ be the set of fully-PPT states, $\mathrm{BS}$ the
biseparable states and $\operatorname{rob}(\rho)=\min\{s\ge0:(\rho+sI/d)/(1+s)\in\mathrm{BS}\}$, so that
$p_{\rm GME}(\rho)=\operatorname{rob}(\rho)/(1+\operatorname{rob}(\rho))$. Then
\begin{equation}\label{eq:ppt3-robust-dual}
 \operatorname{rob}(\rho)=\max\Bigl\{0,\ \max_{W\in\mathrm{BS}^*,\ \tr W=d}-\tr(W\rho)\Bigr\},
\end{equation}
where $\mathrm{BS}^*$ is the cone of Hermitian operators nonnegative on $\mathrm{BS}$. Consequently $\operatorname{rob}$ is convex,
its maximum over $\mathcal P$ is attained at an extreme point of $\mathcal P$, and so is that of the increasing function
$p_{\rm GME}=\operatorname{rob}/(1+\operatorname{rob})$, and for a fixed witness the best state is
the solution of a semidefinite program. The same holds with $\mathrm{BS}$ replaced by any closed convex set of states
containing $\mathrm{BS}$, for example the states passing every cyclic level-two test of Lemma~\ref{lem:ppt3-dps}.
\end{proposition}
\begin{proof}
$I/d$ is interior to the cone generated by $\mathrm{BS}$: it is a full-support average of product pure states, whose
projectors span the Hermitian operators, so no nonzero element of $\mathrm{BS}^*$ vanishes on it. Hence the conic program
defining $\operatorname{rob}$ is strictly feasible, and conic duality gives \eqref{eq:ppt3-robust-dual} with no gap; each
$W\in\mathrm{BS}^*$ with $\tr W=d$ gives $s\ge-\tr(W\rho)$. Since $I/d$ is interior, the slice
$\{W\in\mathrm{BS}^*:\tr W=d\}$ is compact, so $\operatorname{rob}$ is a maximum of affine functions over a compact set and hence
continuous and convex. By Bauer's maximum principle it attains its maximum over the compact convex set $\mathcal P$ at an
extreme point. For fixed $W$,
minimising $\tr(W\rho)$ over $\rho\succeq0$, $\tr\rho=1$, $\rho^{T_S}\succeq0$ is a semidefinite program.
\end{proof}

\begin{remark}[where the robust states sit]\label{rem:ppt3-robust-structure}
\cls{numerical evidence; DS-7 generator, DS-4 generator}
The search for $\tilde\rho$ alternated the two steps of Proposition~\ref{prop:ppt3-noise-extremal}:
optimising the witness for a state and the fully-PPT state for a witness. For states invariant under the cyclic permutation and global parity, the
state is block diagonal over six sectors labelled by parity and cyclic charge (the partial-transpose constraints mix
charges and respect only the two parity blocks): the charge-zero sectors spanned by
$\{\ket{000},\ket{011}+\ket{101}+\ket{110}\}$ and $\{w,v\}$, and the one-dimensional sectors $e,\bar e,o,\bar o$.
The following observations come from floating-point searches and are not certified. Extreme points differ in how their
rank is distributed among these sectors. The surface of Section~\ref{sec:ppt3-surface}
occupies $u_s$, the whole $\{w,v\}$ sector, $e$ and $o$, while the optimiser behind $\tilde\rho$ occupies $u_s$, one
direction $w+cv$, $e$, $\bar e$ and $o$, with $s\approx0.974$ and $c\approx-0.750$: the same symmetries as the surface,
but one combination of $w$ and $v$ is traded for the opposite-chirality vector $\bar e$. Both have profile $(5,5,5,5)$.
The level-two white-noise tolerance of $\rho_{\rm opt}$ before mixing is numerically $2.36\times10^{-2}$; whether some fully-PPT
three-qubit state is more robust is open.
\end{remark}

\section{PPT mixtures and the remaining boundary}\label{sec:ppt3-relaxations}

The examples separate full PPT from biseparability. We explain why PPT-mixture tests miss them, show that the supported
witness method applies to every minimal-rank example, and distinguish a remaining question about maximally mixed marginals.

\subsection{Why PPT-mixture tests cannot detect the states}

A \emph{PPT mixture} is a convex combination of states, each PPT
across one of the three cuts~\cite{JMG,NMG}. Every biseparable
state is a PPT mixture, but the converse fails~\cite{HK}.

\begin{corollary}[PPT mixtures and outer relaxations]\label{cor:ppt3-ppt-mixture}
\cls{proved; the level-two detection is the exact-rational certificates of Section~\ref{sec:ppt3-noise}}
Every fully-PPT GME state in Theorems~\ref{thm:ppt3-family}
and~\ref{thm:ppt3-orbit} and Corollaries~\ref{cor:ppt3-extreme6},
\ref{cor:ppt3-extreme5} and~\ref{cor:ppt3-robust} is a
PPT mixture. No PPT-mixture witness can detect any of them.
The level-two symmetric-extension certificates of
Section~\ref{sec:ppt3-noise} do detect the first family and
the states $\rho_*$, $\sigma$ and $\tilde\rho$ of Corollaries~\ref{cor:ppt3-noise-reused},
\ref{cor:ppt3-extreme5} and~\ref{cor:ppt3-robust}; we do not assert a level-two certificate for
every member of the local-filter surface.
\end{corollary}
\begin{proof}
Every fully-PPT state is itself a one-term PPT mixture. The GME and level-two witness claims
were proved in the preceding sections.
\end{proof}

\begin{remark}[rank-four PPT pieces do not always suffice]\label{rem:ppt3-splitting}
\cls{proved}
The family of Theorem~\ref{thm:ppt3-family} also rules out a proposed universal decomposition
of fully-PPT rank-five states into positive operators
$\rho=\sigma_A+\sigma_B+\sigma_C$ with
$\sigma_c^{T_c}\succeq0$ and $\rank\sigma_c\le4$.
By the $2\times4$ low-rank theorem~\cite{KCKL,HLVC},
each summand would be separable across its indicated cut,
making $\rho$ biseparable.
\end{remark}

\subsection{Every minimal-rank example has the same geometry}\label{sec:ppt3-moment-cone}

The supported-witness method extends beyond our chosen vectors. For a fully-PPT GME state of minimal rank, the kernel contains no cut-product
vector; this forces the product rays at each cut to form a rational quartic curve spanning the support.
The three curves determine all its biseparable states and all supported witnesses.

\begin{theorem}[three images of one quartic moment cone]\label{thm:ppt3-universal-moments}
\cls{proved}
Let $R\subset(\mathbb C^2)^{\otimes3}$ be a five-dimensional subspace whose orthogonal complement contains no nonzero vector
that is a product across any of the three cuts; by Lemma~\ref{lem:ppt3-product-kernel} (Appendix~\ref{sec:ppt3-app-moment}) this holds for the range of every
rank-five fully-PPT GME state. Let $B$ be an $8\times5$ basis matrix of $R$, put $\nu(a)=(a_0^4,a_0^3a_1,a_0^2a_1^2,
a_0a_1^3,a_1^4)^T$ for $a\in\mathbb C^2$, and let $\mathcal M_4=\operatorname{cone}\{\nu(a)\nu(a)\dg\}$ be the spin-two
(quartic) moment cone. For each cut $c$ there is an invertible $5\times5$ matrix $T_c$ such that the cut-$c$ product
rays in $R$ are exactly the rays of $BT_c\nu(a)$, $a\neq0$. Consequently, for Hermitian $X$,
\begin{equation}\label{eq:ppt3-universal-bs}
 BXB\dg\ \text{is biseparable}\iff X\in\sum_{c=A,B,C}T_c\mathcal M_4T_c\dg ,
\end{equation}
every biseparable operator on $R$ is a sum of at most $25$ cut-product projectors, and a Hermitian $5\times5$ matrix $H$
satisfies $\tr(HX)\ge0$ for every biseparable $BXB\dg$ if and only if
\begin{equation}\label{eq:ppt3-universal-dual}
 \nu(a)\dg T_c\dg HT_c\,\nu(a)\ge0\qquad\text{for all }a\in\mathbb C^2,\ c=A,B,C .
\end{equation}
\end{theorem}

\emph{Idea of the proof.} A fully-PPT state whose kernel contains a product vector is separable across that cut, so the
kernel of a minimal-rank example is product-free. The product vectors $a\otimes b$ in $R$ are then the solutions of a
$3\times4$ linear pencil in $a$ of constant rank three, which is forced to be a single Kronecker block; its kernel is
cubic in $a$, so the product vectors trace a rational normal quartic spanning $R$. Appendix~\ref{sec:ppt3-app-moment}.

Thus, on the support of a minimal-rank example, biseparability is decided by where three copies of one universal cone
sit relative to each other, and genuine multipartite entanglement is certified exactly by a $5\times5$ Hermitian $H$
whose three quartic forms~\eqref{eq:ppt3-universal-dual} are nonnegative while $\tr(HX)<0$. Such an $H$ exists for
every GME state supported on $R$, since the cone in~\eqref{eq:ppt3-universal-bs} is closed. The witnesses $H_s$ and $H_0$ of
Sections~\ref{sec:ppt3-family} and~\ref{sec:ppt3-surface} are instances, with the quartic forms factorised by hand.
Testing the nonnegativity of a quartic form in one complex variable is a small problem: a real polynomial of degree eight
in two variables.

\subsection{The flat-marginal question}\label{sec:ppt3-flat}

A three-qubit state is \emph{flat} if every proper marginal
is maximally mixed. Flat states are a natural place to look for further examples, because every flat state is automatically fully
PPT (shown below), so only biseparability is in question. Our examples settle existence of fully-PPT GME, but leave the following
conjecture open:
\begin{conjecture}[flat rank-five states]\label{conj:ppt3-flat}
\cls{open; repository name Conjecture R}
Every flat three-qubit state of rank five is biseparable.

\end{conjecture}
Every flat three-qubit state is fully PPT, with rank profile $(r,r,r,r)$, $r=\rank\rho$. Indeed
$\rho=I/8+T$ with $T$ a real combination of three-body Pauli products. Partial transposition on qubit $c$ followed by
conjugation with $\sigma_y$ on that qubit negates every Pauli factor on $c$, and the antiunitary
$J=\sigma_y^{\otimes3}K$ ($K$ complex conjugation) negates all three factors, so
\[
 \sigma_y^{(c)}\rho^{T_c}\sigma_y^{(c)}=I/8-T=J\rho J^{-1}.
\]
Each partial transpose is therefore isospectral to $\rho$. Profiles are unchanged by invertible local filtering.

\begin{remark}[rank profile is not flatness]\label{rem:ppt3-flat-profile}
\cls{proved}
The first family has profile $(5,8,8,8)$ in the PPT interval
and $(5,7,7,7)$ at its endpoints, while
Corollary~\ref{cor:ppt3-extreme6} has profile $(5,6,6,6)$.
None is locally filter-equivalent to a flat rank-five state.
The second family and its extreme member \emph{do} have profile
$(5,5,5,5)$, yet they are GME. Thus that profile by itself
cannot imply biseparability; Conjecture~\ref{conj:ppt3-flat} must use the
flat-marginal condition, not only the ranks.
\end{remark}

\subsection{Relation to earlier work}\label{sec:ppt3-prior}

Fully-PPT genuine multipartite entanglement is known in systems larger than three qubits. For states supported on the totally
symmetric subspace, any cut-product vector in the support
is fully product, so biseparability equals full separability
there. This connects the symmetric multipartite PPT-entangled
constructions for four or more qubits~\cite{TG,TAHKSL,ATSL} to GME;
mere permutation invariance, without support in
the symmetric subspace, is not sufficient for this inference.
Fully-PPT GME examples also exist for three qutrits
and higher local dimensions~\cite{HS,PM,SB,SBZL}; for~\cite{SBZL} the state is the normalised projector onto the
complement of a genuinely unextendible product basis.

For three qubits, Ac\'in, Bru\ss, Lewenstein and Sanpera gave an early construction of fully-PPT three-qubit edge states, with rank
profiles up to $(7,7,7,7)$, including an X-shaped family~\cite{ABLS}, of which Kampermann et al.\ generated a pseudo-bound-entangled NMR version~\cite{KPBS}.
By Han and Kye, every X-shaped fully-PPT multi-qubit state is biseparable~\cite[Cor.~5.4]{HanK}; these NMR
states are therefore not GME, and, since biseparability and PPT are invariant under local unitaries, the GME states of this
paper are not X-shaped in any product basis. Explicit three-qubit PPT entangled states of low rank are also given
in~\cite{HaKye15}.
A PPT state is an edge state across a cut if no product vector of its range has its partial conjugate in the range of the partial
transpose. The states of the $(5,5,5,5)$ surface are edge states across every cut (Corollary~\ref{cor:ppt3-extreme5}). The interior states $\rho_t$ are not: their partial transposes are invertible
and the range contains the product vector $\ket{111}$.

Ha and Kye gave a three-qubit PPT mixture that is not biseparable, but their rank-seven example was established
PPT across a selected cut, not all three~\cite{HK}. This already refuted the conjecture of Novo, Moroder and G\"uhne that
PPT mixtures coincide with biseparable states for three qubits~\cite{NMG}. The states here are a stronger counterexample: a
single-term PPT mixture that is PPT across all three cuts. By Observation~3 of~\cite{NMG}, a permutation-invariant
three-qubit state is biseparable if and only if it is a PPT mixture, so no state of this paper is permutation invariant, in
line with Corollary~\ref{cor:ppt3-symmetrisation}.
Garberg, Irgens and Myrheim found fully-PPT states
numerically with many rank profiles, including $(5,5,5,5)$ and extreme states, without a GME
certificate~\cite{GIM}. Their phrase ``genuine tripartite entanglement'' for rank-$(4,4,4,4)$ states means entangled while
separable across every bipartition, which is not GME in the present sense.
Hari Krishnan, Ranjan and Banik~\cite{KRB} provided
fully-PPT three-qubit states outside
$\bigcap_c\mathrm{Sep}_c$, the intersection of the
three cut-separable sets, by cyclically symmetrising an embedded $2\times4$ PPT state and subtracting low-rank terms.
This does not imply GME:
the biseparable set is the larger convex hull
$\mathrm{BS}=\operatorname{conv}(\bigcup_c\mathrm{Sep}_c)$.
Our supported witness excludes the \emph{convex hull},
not merely the intersection. T\'oth and V\'ertesi found numerically a three-qubit state that is PPT across all
bipartitions and nevertheless useful for metrology~\cite{TV}; it is entangled but was not shown to be GME.
Singh and Banik noted that no explicit three-qubit state that is PPT across all three cuts and GME was known,
and showed that, if one exists, it can be taken extreme with all four ranks at least five and the smallest equal to five
or six~\cite{SB}; they did not settle existence. Our extreme states satisfy these restrictions.

To our knowledge these are the first explicit
\emph{three-qubit} examples proved both fully PPT and GME;
the claim does not extend to more parties
or higher local dimensions.



\section{Discussion}\label{sec:ppt3-discussion}

Physically, the examples separate two questions that are easy to conflate: whether entanglement can be extracted across
some bipartition, and whether it is shared by all three parties. Even for three qubits, the first does not decide the second. A
fully-PPT GME state yields no singlet across any cut and no GHZ state under any LOCC protocol, yet it is not a mixture of
states each separable across some bipartition. The relevant gap is between the intersection of the three PPT sets and
the convex hull of the three cut-separable sets. The support of $\rho_t$ contains product vectors across every cut, but no convex mixture of them
reproduces $\rho_t$; the supported witness certifies this failure of mixing. Symmetry plays a definite role.
Reflection symmetry together with PPT forces separability across a cut (Corollary~\ref{cor:ppt3-symmetrisation}), so the
states must break every exchange of two qubits; ours do so through the cyclic phases of the $W$-type vectors, and they
are also locally chiral (Corollary~\ref{cor:ppt3-chirality}).

For certification, PPT-mixture tests cannot certify the genuine entanglement of these states; a stronger relaxation of
biseparability is one way to do so. The second level of the
symmetric-extension hierarchy suffices here, and for white noise its certified margins nearly match explicit biseparable
decompositions (Proposition~\ref{prop:ppt3-noise-sharp}), so the small size of the margins is a property of the states,
not of the relaxation. By
Proposition~\ref{prop:ppt3-noise-extremal} a most robust state can be taken extreme in the fully-PPT set, and numerical searches
suggest that within one symmetry class the extreme points differ in how their rank is distributed among the symmetry
sectors (Remark~\ref{rem:ppt3-robust-structure}). An NMR experiment has generated pseudo-bound-entangled fully-PPT three-qubit
states that are not GME~\cite{KPBS}. The state $\tilde\rho$ remains GME under depolarising errors of up to one per cent
per qubit, but generic preparation errors need much tighter control: its certified PPT-stability radius, $2.0\times10^{-5}$, is far
smaller than its GME radius, $7.3\times10^{-3}$ (Corollary~\ref{cor:ppt3-robust}); certifying such a state also requires witnesses that go beyond
partial transposition.

Several questions remain open: whether some fully-PPT three-qubit state tolerates substantially more noise, how to
describe the family of $\tilde\rho$ analytically, and how to decompose a witness into practical local measurements. The mixed-$W$ conjecture of
Ac\'in et al.~\cite{ABLS} holds for all our explicit states, including $\tilde\rho$ (Theorem~\ref{thm:ppt3-w-class}), but
remains open for general PPT entangled three-qubit states. A separate problem is whether every rank-five three-qubit state with maximally mixed one- and
two-qubit marginals is biseparable (Conjecture~\ref{conj:ppt3-flat}). Our rank-$(5,5,5,5)$ examples show that
the rank profile alone cannot prove it; the flat-marginal condition must play an essential role.

\section*{Use of AI}
All results, documents and scripts were generated by AI agents (mainly Claude Opus and Sonnet 5.5, and OpenAI's GPT 6.1 sol)
with human supervision and many cross-checks, including AI-generated Lean verification of the main analytic results.

\section*{Acknowledgements}
We acknowledge support by the UK STFC grant ST/T000473/1.

\appendix
\small
\section{Computations for the rank-five family}\label{sec:ppt3-app-family}

This appendix contains the proofs of the statements of Section~\ref{sec:ppt3-family}, whose notation it uses.

\begin{proof}[Proof of Theorem~\ref{thm:ppt3-family}, PPT and rank claims]
Cyclic permutations of the qubits leave $\rho_t$ invariant, so
its three partial transposes are unitarily equivalent. The
unnormalised numerator of $\rho_t^{T_A}$ commutes with the global parity $Z^{\otimes3}$, so it splits into even and
odd parity blocks. With $\beta=3+\omega$ and $\gamma=t+3\omega^2$ ($|\beta|^2=7$, $|\gamma|^2=t^2-3t+9$), in the respective basis orders
$(000,011,101,110)$ and $(001,010,100,111)$ these blocks of $7(t+3)\rho_t^{T_A}$ are
\[
\text{even: }\begin{pmatrix}
t & -t & \beta & \bar\beta\\
-t & t+3 & 0 & 0\\
\bar\beta & 0 & t+3 & \gamma\\
\beta & 0 & \bar\gamma & t+3
\end{pmatrix},\qquad
\text{odd: }\begin{pmatrix}
4 & \beta & -t & \bar\gamma\\
\bar\beta & 4 & -t & \gamma\\
-t & -t & 4 & 0\\
\gamma & \bar\gamma & 0 & 3t
\end{pmatrix},
\]
and their leading principal minors are
\begin{equation}\label{eq:ppt3-minors}
\begin{aligned}
\text{even: }&
t,\quad 3t,\quad (t+3)(3t-7),\quad
3(8t^2-15t-36),\\
\text{odd: }&
4,\quad 9,\quad 3(12-t^2),\quad
-3(t-3)(3t^2+4t-48).
\end{aligned}
\end{equation}
The even determinant is nonnegative exactly when $t\ge t_-$.
For $3\le t<2\sqrt3$, the first three leading minors
of each block are positive, while the odd determinant is
nonnegative exactly when $t\le t_+$. Since
$3<t_-<10/3<t_+<2\sqrt3$, the leading $3\times3$ corners are positive definite throughout $[t_-,t_+]$.
Each block is positive semidefinite exactly when its scalar Schur complement, or equivalently its determinant, is
nonnegative. This gives PPT precisely on $[t_-,t_+]$.
Outside this interval one of the displayed determinants
is negative. At either endpoint exactly one block has
rank three (its leading $3\times3$ minor stays positive);
the other has rank four. The stated profiles follow.
The GME claim is established in Section~\ref{sec:ppt3-witness}, after constructing
a witness on the range.
\end{proof}

\begin{proof}[Proof of Proposition~\ref{prop:ppt3-witness}]
\emph{(a).} The product condition reduces to four linear equations. Write a general vector of $R_s$ as
$B_s(U,V,W,E,O)^T$. It is a product with the first-qubit ray
$(1,z)^T$ precisely when the two $BC$ rows are proportional
with ratio $z$. In the computational basis this is
\[
\begin{aligned}
-zU+W+O&=0,&
-sU-zW+\omega E-z\omega^2O&=0,\\
-sU-zW+\omega^2E-z\omega O&=0,&
szU+V-zE&=0.
\end{aligned}
\]
Subtracting the middle equations and using
$\omega+\omega^2=-1$ gives
\[
 E=-zO,\qquad W=zU-O,\qquad
 (s+z^2)U=2zO,\qquad V=-szU-z^2O.
\]
For every finite $z$, the solutions form a one-dimensional
space, because $s+z^2$ and $2z$ never vanish simultaneously
when $s>0$. Setting $U=2z$, $O=s+z^2$ gives exactly
$q_s(z)$ in \eqref{eq:ppt3-curve}; it remains nonzero at
$z=0$ and $z^2=-s$. At the projective point
$(0,1)^T$ (i.e.\ $z=\infty$), the upper $BC$ row must
vanish, which forces $U=W=E=O=0$ and leaves only $V$,
the ray of $v$. This proves the full cut-product classification.

For Appendix~\ref{sec:ppt3-kernel}, $R^\perp$ is spanned by the three symmetry-sector vectors omitted from $R$, namely
$\bar e$, $\bar o$ and $3\ket{000}+\ket{011}+\ket{101}+\ket{110}$: the $3\times8$ matrix $K$ whose rows are their bras,
\[
K_0=[0,0,0,1,0,\omega,\omega^2,0],\quad
K_1=[0,\omega^2,\omega,0,1,0,0,0],\quad
K_2=[3,0,0,1,0,1,1,0]
\]
has rank three and $KB=0$, so $R=\ker K$.

\emph{(b).} A rescaling removes the parameter $s$: with
$D_s=\operatorname{diag}(1,s^{3/2},s^{1/2},s,s^{1/2})$,
\[
 q_s(\sqrt s\,z)=\sqrt s\,D_sq_1(z),\qquad
 H_s=s^3D_s^{-1}H_5D_s^{-1},
\]
so $q_s(\sqrt s z)\dg H_sq_s(\sqrt s z)=s^4\,q_1(z)\dg H_5q_1(z)$, and the right side of \eqref{eq:ppt3-sos} at
$(s,\sqrt s z)$ is $s^4$ times its value at $(1,z)$. It suffices to take $s=1$, which is \eqref{eq:ppt3-h5}. Put $d=\operatorname{Re}(z^2)=x^2-y^2$, so
$r+d=2x^2$. The squared moduli of the entries of $q_1(z)$ are
\[
4r,\quad r^2(r^2+9+6d),\quad r^2+1-2d,\quad r(r^2+1+2d),\quad r^2+1+2d,
\]
and $H_5$ is diagonal, so
\begin{align*}
q_1\dg H_5q_1&=-48r+r^2(r^2+9+6d)+11(r^2+1-2d)-6r(r^2+1+2d)+14(r^2+1+2d)\\
&=r^4-6r^3+34r^2-54r+25+6d(r-1)^2 .
\end{align*}
Since $(r-1)^2(r^2-4r+25)=r^4-6r^3+34r^2-54r+25$ and $12x^2=6(r+d)$, this equals
$(r-1)^2\bigl[(r-5)^2+12x^2\bigr]$. This proves \eqref{eq:ppt3-sos}.
The cyclic permutation of qubits fixes $u_s,v,w$ and
acts on $e,o$ by phases, which commute with the diagonal $H_s$, so the form controls the other cuts as well.
The point at infinity contributes $H_{vv}=1$. Finally, every positive pure-state decomposition of an operator
supported on $R_s$ uses vectors in $R_s$: its expectation
on $R_s^\perp$ vanishes, and each positive summand must
vanish there separately. Thus the nonnegative curve
values also control mixtures across different cuts.
\end{proof}

\begin{proof}[Proof of Corollary~\ref{cor:ppt3-chirality}]
\emph{(a).} The five nonzero eigenvalues of $\rho_t$ are, up to the common
factor $1/[7(t+3)]$, $4t,3t,9,9,3$, on the respective
rays $(u,v,w,e,o)$. Since $t>3$ throughout
\eqref{eq:ppt3-interval}, the two real eigenrays
$u$ and $v=\ket{111}$ are each nondegenerate.
Because $u$ and $v$ are real, $\overline{\rho_t}$ has the same
nondegenerate eigenrays, so a product unitary mapping $\rho_t$ to
its complex conjugate must fix the ray of $v=\ket{111}$. Each
factor then fixes the ray of $\ket1$, hence also the orthogonal
ray of $\ket0$, and is diagonal. Write its relative local phases as
$(\theta_A,\theta_B,\theta_C)$. Fixing the $u$ ray forces
\[
\theta_A+\theta_B\equiv\theta_A+\theta_C
\equiv\theta_B+\theta_C\equiv0\pmod{2\pi},
\]
so all three phases are $0$ or all are $\pi$.
Up to a scalar, the only candidates are $I$ and
$Z\otimes Z\otimes Z$, both of which fix $\rho_t$.
But, in the computational basis,
\[
  \langle011|\rho_t|101\rangle
       =\frac{t+3\omega^2}{7(t+3)}
       \notin\mathbb R .
\]
Thus $\rho_t\ne\overline{\rho_t}$, and neither candidate
implements complex conjugation.

\emph{(b).} By Proposition~\ref{prop:ppt3-witness}(a) the $A|BC$-product rays of $R_s$ are $[v]$ and
$[B_sq_s(z)]=[(\ket0+z\ket1)\otimes\ket{\phi_s(z)}]$, $z\in\mathbb C$, and substituting \eqref{eq:ppt3-curve} shows that the $2\times2$
coefficient matrix of $\ket{\phi_s(z)}$ has determinant $-3(z^2+s)^2$. So $\ket{\phi_s(z)}$ is a $B|C$ product exactly when
$z=\pm i\sqrt s$, where $B_sq_s(\pm i\sqrt s)=\pm2i\sqrt s\,(\ket0\pm i\sqrt s\ket1)^{\otimes3}$. Hence $R_s$ contains exactly three fully
product rays, $P_\infty=[\ket1^{\otimes3}]$ and $P_\pm=[(\ket0\pm i\sqrt s\ket1)^{\otimes3}]$; they span $\operatorname{span}(u_s,v,w)$, because
$(\ket0\pm i\sqrt s\ket1)^{\otimes3}=u_s\pm i\sqrt s\,(w-sv)$. The conjugate $\overline{R_s}$ contains the conjugate rays, the same three with $P_+$ and $P_-$ exchanged. An invertible
local filter $F$ that maps $R_s$ onto $R_s$ or onto $\overline{R_s}$ permutes the three rays. Each ray is $[\chi^{\otimes3}]$, so $F_A$,
$F_B$ and $F_C$ all induce the same permutation of the three distinct points $[\chi]\in\mathbb{CP}^1$, and since a M\"obius map is determined by its action on three
points, $F=cF_0^{\otimes3}$. Such an $F$ commutes with the cyclic shift $\Pi$, so it maps the $\omega^2$-eigenspace of $\Pi$ into itself. That eigenspace is
$\operatorname{span}(e,o)\subset R_s$, whereas $\overline{R_s}$ meets it in $0$: $\overline{R_s}\ni\bar e,\bar o$ lies in the $\omega$-eigenspace and the other three
generators of $\overline{R_s}$ are $\Pi$-invariant. So $F(R_s)\ne\overline{R_s}$. If a state $\rho$ with range $R_s$ satisfied $\rho=F\bar\rho F\dg$ (or the same with $Y^{\otimes3}\bar\rho Y^{\otimes3}$, a local-unitary image of $\bar\rho$), then
there would be an invertible local filter $F'$ ($F^{-1}$, respectively $Y^{\otimes3\dagger}F^{-1}$) with $F'(R_s)=\overline{R_s}$, which the argument excludes. For the stabiliser, each of the six permutations of the three points is
induced by a unique M\"obius map, and $F_\pi^{\otimes3}$ permutes $P_\infty,P_\pm$, hence preserves $\operatorname{span}(u_s,v,w)$, and as it commutes with $\Pi$ it preserves
the $\omega^2$-eigenspace $\operatorname{span}(e,o)$ of $\Pi$ (which is all of that eigenspace in $(\mathbb C^2)^{\otimes3}$). So the stabiliser is $\{[F_\pi]^{\otimes3}\}\cong S_3$.

\end{proof}

\begin{proof}[Proof of Corollary~\ref{cor:ppt3-extreme6}]
Write the numerator of $\sigma_6$ as $a\,uu\dg+b\,vv\dg+3\,ww\dg+3\,ee\dg+oo\dg$; it is the numerator of
$\rho_a$ with the weight $3a$ of $vv\dg$ replaced by $b$. The vector $v=\ket{111}$ enters $\sigma_6^{T_A}$ only through the
$(111,111)$ entry, the bottom diagonal entry of the odd block, and equals $b$ there. Hence the even block, and the three leading
principal minors of the odd block, are those of $\rho_a^{T_A}$ in \eqref{eq:ppt3-minors}: at $a=t_-$ the even minors
$a,\ 3a,\ (a+3)(3a-7)$ are positive and, since $8a^2-15a-36=0$, the even determinant vanishes, so the even block is positive
semidefinite of rank three. The odd minors $4,\ 9,\ 3(12-a^2)$ are positive ($t_-<2\sqrt3$), and the odd determinant is
affine in $b$,
\[
 \det(\text{odd block})=3(12-a^2)\,b+3(5a^2+24a-144),
\]
which equals $-3(t-3)(3t^2+4t-48)$ of \eqref{eq:ppt3-minors} at $a=t$, $b=3t$. It vanishes at
$b=(5a^2+24a-144)/(a^2-12)$, and this equals $8(9-a)/5$ modulo $8a^2-15a-36=0$, since
$5(5a^2+24a-144)-8(9-a)(a^2-12)=8a^3-47a^2+24a+144$ is divisible by $8a^2-15a-36$
(the quotient is $a-4$ and the remainder is $0$).
Hence at $b=8(9-a)/5$ the odd block is positive semidefinite of rank three as well. Thus each partial transpose of $\sigma_6$ has
rank six.
The supported witness has expectation
\[
\tr(H_5X_6)
=\frac{b-12a+29}{4a+b+21}
=\frac{217-68a}{5(4a+b+21)}<0,
\]
because $a>13/4>217/68$. This proves GME by
Proposition~\ref{prop:ppt3-witness}.
For extremality, $\sigma_6$ is extreme in the set of fully-PPT states exactly when the only Hermitian $D$ with range in $\operatorname{ran}\sigma_6$
and $D^{T_c}k=0$ for every kernel vector $k$ of $\sigma_6^{T_c}$, at each cut $c$, is a multiple of $\sigma_6$ (the supported-face
criterion of~\cite{GIM}). Write $D=BX_DB\dg$ with $X_D$ Hermitian. The conditions
$D^{T_c}k=0$ are a real-linear system on the $25$-dimensional real space of such $X_D$. To compute its rank exactly, complexify:
let $X_D$ range over all complex $5\times5$ matrices (25 complex unknowns) and impose both $D^{T_c}k=0$ and $k\dg D^{T_c}=0$ for the
kernel vectors $k$, at each cut. The two sets are adjoint to each other for Hermitian $D$, so the solution space is closed under
$X_D\mapsto X_D\dg$, and is the complexification of the real solution space on Hermitian $X_D$ (split $X_D$ into its
Hermitian and anti-Hermitian parts). Real nullity therefore equals complex nullity. The complexified system has rank $24$ over
$\mathbb Q(\sqrt{17},i\sqrt3)$, which contains the entries of $\sigma_6$ and of the kernel vectors, so its nullspace, and hence the
real nullspace, is exactly the scalar direction of \eqref{eq:ppt3-sigma6}.
\end{proof}

\section{\texorpdfstring{Computations for the profile-$(5,5,5,5)$ surface}{Computations for the profile-(5,5,5,5) surface}}\label{sec:ppt3-app-surface}

This appendix contains the proofs of the statements of Section~\ref{sec:ppt3-surface}, whose notation it uses.

\begin{proof}[Proof of Lemma~\ref{lem:ppt3-h0}]
At fixed radius, the two-variable identity reduces to quadratic interpolation. Hold $r=\xi+\zeta>0$ fixed, put
$P(\xi,r-\xi)=q_1^\dagger H_0q_1$, and substitute $z=i\sqrt r$
and $z=\sqrt r$ into the five entries of $q_1$.
Using $x_0^2q=y_1^2$ and collecting powers gives the axis values
\[
 P(0,r)=(r-1)^2(r-y_1)^2,\qquad
 P(r,0)=(r-x_0)^2(r^2+pr+q).
\]
For fixed $r$, each diagonal term of $P$ is affine in
$\xi$ (it depends on $z^2$ only through
$\operatorname{Re}(z^2)=2\xi-r$).
The off-diagonal $vw$ entry contributes a quadratic
whose $\xi^2$ coefficient is $16(H_0)_{vw}=-16\nu$.
Thus $P(\xi,r-\xi)$ is the unique quadratic with the displayed
endpoint values and that leading coefficient, namely
\[
 P(\xi,\zeta)=\frac{\zeta}{r}P(0,r)+\frac{\xi}{r}P(r,0)
       +16\nu\,\xi\zeta,\qquad \xi+\zeta=r>0.
\]
Multiplying by $r$ gives \eqref{eq:ppt3-h0}. All terms are nonnegative since $p,q,\nu,\xi,\zeta\ge0$,
so $P\ge0$ for $r>0$ and also at $r=0$ by continuity.
At the point $z=\infty$, the product ray is $v$ and
$H_{vv}=1$. Cyclic invariance holds because the only off-diagonal entries of $H_0$
couple $v$ and $w$, both fixed by the cyclic permutation; so the form controls
all three cuts and their biseparable mixtures.
\end{proof}

\begin{proof}[Proof of Theorem~\ref{thm:ppt3-orbit}]
\emph{Filter reduction.} The formulas for $F_\lambda$ follow
by counting computational-basis excitations. Under
$aF_\lambda(\,\cdot\,)F_\lambda^\dagger$ with
$\lambda^2=1/a$, the coefficients $(\gamma,\beta,V,G,E,O)$ of $u_su_s\dg,ww\dg,\dots$ at
$\gamma=\alpha$ become
$(a\alpha,\beta,V/a^2,G/a,E/a,O)$, exactly as in
\eqref{eq:ppt3-weights}. It therefore suffices to prove the remaining assertions at $\alpha=\gamma=1$.

\emph{Rank and PPT.}
Put $\kappa=1-\tau^2$. Here $\rho_\tau=\rho_{\tau,1}=BX_\tau B\dg$ is the unnormalised surface operator and $X_\tau$ its
coefficient matrix; they are not the state $\rho_t$ and matrix $X_t$ of Section~\ref{sec:ppt3-family}. Introduce the positive
coordinate
\[
Q=\tau+\sqrt{1+\tau^2},\qquad
\tau=\frac{Q-Q^{-1}}{2},\quad
Q^2+Q^{-2}=2(1+2\tau^2),\quad
1<Q<1+\sqrt2 .
\]
The weights simplify to
\[
O=\kappa/\tau,\quad E=\kappa Q^2,\quad G=(E-1)/\tau,\quad
V=(E^2-2\kappa E+\kappa)/(\tau\kappa),
\]
and in terms of $Q$ alone $\kappa=(6Q^2-Q^4-1)/(4Q^2)$, $E=(6Q^2-Q^4-1)/4$ and $G=Q(5-Q^2)/2$, so every identity below is a
direct substitution.
The vectors $(u,v,w,e,o)$ are linearly independent, and the only off-diagonal coefficient of $\rho_\tau$ is its
$w$--$v$ coupling. Since $\kappa,E,O,\tau>0$,
positivity is equivalent to
\begin{equation}\label{eq:ppt3-surface-simpledet}
\tau V-G^2=2Q^2\kappa(Q^2-3)\ge0.
\end{equation}
Equality gives rank four and strict inequality rank five.

There is also an explicit Gram representation of the
\emph{odd} parity block of $\rho_\tau^{T_A}$, in the order
$(001,010,100,111)$. Define
\[
L=\begin{pmatrix}\tau+O&\tau+\omega O\\
\tau+\omega^2O&\tau+O\end{pmatrix},
\qquad
M=\begin{pmatrix}-1&1+\omega E\\
                 -1&1+\omega^2E\end{pmatrix}.
\]
Here $\det L=3\tau O=3\kappa>0$, and direct multiplication
using $\tau(\tau+O)=1$ gives
\[
M^\dagger L^{-1}M=
\begin{pmatrix}\tau+O&G\\G&V\end{pmatrix},
\qquad
(\rho_\tau^{T_A})_{\rm odd}=
\begin{pmatrix}L&M\\M^\dagger&M^\dagger L^{-1}M\end{pmatrix}.
\]
Hence the odd block is positive semidefinite of rank two,
for every $0<\tau<1$. The \emph{even} block, in the order
$(000,011,101,110)$, is
\[
\begin{pmatrix}
1&-1&b&\bar b\\
-1&1+E&G&G\\
\bar b&G&1+E&1+\omega^2E\\
b&G&1+\omega E&1+E
\end{pmatrix},\qquad b=\tau+\omega O.
\]
Its determinant factors as
\[
\frac{3(\kappa-\tau^2E)}{\tau^4}
 \bigl[E^2-2\kappa(1+2\tau^2)E+\kappa^2\bigr]=0,
\]
the last equality following from $E=\kappa Q^2$ and
$Q^2+Q^{-2}=2(1+2\tau^2)$.
Its first three leading principal minors are
\begin{equation}\label{eq:ppt3-surface-minors}
1,\qquad E,\qquad
\frac{(Q^2-3)(Q^4+3)\kappa^2}{4\tau^2}.
\end{equation}
For $Q^2>3$ the leading $3\times3$ corner is positive
definite; the full determinant is zero, so the Schur
complement vanishes and this block is positive semidefinite
of rank three. At $Q^2=3$ positivity follows by
continuity. Since $Q$ increases with $\tau$ and
$Q^2=3$ precisely at $\tau=1/\sqrt3$, this proves the
sharp positivity window, the rank profile and the
rank-four lower endpoint. Cyclic symmetry gives the
other cuts. The filter reduction proves the conclusions
for every $(\alpha,\beta,\gamma)$, including the rank-at-most-five
assertion as well. The chosen $E$ is the larger root
$\kappa(\tau+\sqrt{1+\tau^2})^2$ of the even-block determinant
factor, and $G,V$ complete the odd-block Gram matrix.

\emph{Genuine entanglement.} At $a=1$ and $3/4\le \tau<1$ (so $2\le Q<1+\sqrt2$; $\tau=3/4$ is $Q=2$), the fixed witness
$H_5$ of Proposition~\ref{prop:ppt3-witness} gives, by substitution of the simplified weights,
\[
 \tr(H_5X_\tau)=-\frac{(Q^2-2Q-1)\,S_6(Q)}{2Q(Q^2-1)},\qquad
 S_6(Q)=Q^6-Q^5-9Q^4+2Q^3+11Q^2+27Q-3 .
\]
With $x=Q-2\in[0,\sqrt2-1)$,
\[
S_6(2+x)=-1-81x-33x^2+50x^3+41x^4+11x^5+x^6<0 .
\]
Indeed $\sqrt2-1<1/2$, and for $0\le x\le1/2$ one has $50x^3\le25x^2$ and
$41x^4+11x^5+x^6\le\bigl(\tfrac{41}8+\tfrac{11}{16}+\tfrac1{32}\bigr)x\le6x$.
Thus $S_6(2+x)\le-1-75x-8x^2<0$, including at $x=0$. Also $Q^2-2Q-1<0$ for $Q<1+\sqrt2$ and
$Q^2-1>0$. Hence $\tr(H_5X_\tau)<0$ throughout $3/4\le \tau<1$.

For the complementary interval $1/\sqrt3<\tau\le3/4$,
so that $\sqrt3<Q\le2$, Lemma~\ref{lem:ppt3-h0}
provides a \emph{fixed} form. Define
\[
P(Q)=Q^4+(8+2\sqrt3)Q^3-(2+4\sqrt3)Q^2
 -(20+14\sqrt3)Q+1+4\sqrt3.
\]
A direct substitution into $\tr(H_0X_\tau)$ factors as
\[
 \tr(H_0X_\tau)=
 -\frac{(Q^2-3)(Q^2-2Q-1)P(Q)}
 {2Q(Q^2-1)}.
\]
Here $P''(Q)>0$ for $Q\ge\sqrt3$,
$P'(\sqrt3)=28+12\sqrt3>0$, and
$P(2)=33-24\sqrt3<0$. Thus
$P(Q)\le P(2)<0$ on $[\sqrt3,2]$.
Since $Q^2-3>0$ and $Q^2-2Q-1<0$,
the trace is strictly negative on
$(1/\sqrt3,3/4]$ and zero at the lower endpoint.
The two witness ranges meet at $\tau=3/4$ and together cover the entire open PPT window, so $\rho_{\tau,1}$ is GME there. Finally, by
Theorem~\ref{thm:ppt3-orbit}(b), $\rho_{\tau,a}=aF_\lambda\rho_{\tau,1}F_\lambda^\dagger$ with $F_\lambda$ invertible and local.
Such a filter, and its inverse, maps a vector that is a product across a cut to a vector that is a product across the same cut,
so it maps biseparable states to biseparable states in both directions, and likewise preserves positivity of every partial
transpose and all ranks. Hence $\rho_{\tau,a}$ is biseparable if and only if $\rho_{\tau,1}$ is, and GME for every $a>0$
follows from the case $a=1$.
\end{proof}

\begin{proof}[Proof of the computations in Remark~\ref{rem:ppt3-endpoint}]
Substituting $q_1(\pm i\sqrt{y_1})$ and $q_1(\pm\sqrt{x_0})$
into~\eqref{eq:ppt3-lower-endpoint} gives $X_{1/\sqrt3}$ directly. A form on $\mathbb C^5$ invariant under the cyclic
($\operatorname{diag}(1,1,1,\omega^2,\omega^2)$), parity and $K\circ S_{BC}$ (complex conjugation) actions on the coordinates
is real and block diagonal on $u$, $(v,w)$, $e$, $o$, so it has six real parameters $(h_u,h_v,h_w,h_e,h_o,h_{vw})$. Vanishing
with vanishing derivative in $r=|z|^2$ at $z=i$, $z=i\sqrt{y_1}$ and $z=\sqrt{x_0}$ (two equations each; the points $-z$ give the
same equations) is a homogeneous linear system of rank five, whose kernel is spanned by $H_0$. Finally
$P(11/5)=(34956-15540\sqrt3)/625>0$, so $\tr(H_0X_{48/55})>0$.
\end{proof}

\begin{proof}[Proof of Corollary~\ref{cor:ppt3-extreme5}]
The parameter $\tau=48/55$ lies in the window of
Theorem~\ref{thm:ppt3-orbit}, which proves positivity and all
four ranks. The displayed coefficients give $\tr N=16323219$;
cyclic and parity invariance follow from the basis vectors.

For an explicit analytic GME certificate, take
$H=H_{1/2}=\operatorname{diag}(-\tfrac32,1,\tfrac{11}4,-3,\tfrac72)$ from Proposition~\ref{prop:ppt3-witness}.
It is nonnegative on every cut-product ray in the range
by \eqref{eq:ppt3-sos}, whereas
\[
\tr(HX_\sigma)=-\frac{493577}{32646438}<0,\qquad
\sigma=B_{1/2}X_\sigma B_{1/2}^\dagger ,
\]
where, in the basis $(u_{1/2},v,w,e,o)$ (note $\tilde u=2u_{1/2}$), $16323219\,X_\sigma$ has
diagonal
\[
(4\cdot825000,\ 2021184,\ 1440000,\ 951720,\ 450625)
\]
and $vw$ entry $145200$, by \eqref{eq:ppt3-sigma5}.

The level-two witness $Y_*$ of Corollary~\ref{cor:ppt3-noise-reused}, defined later in Section~\ref{sec:ppt3-noise}
(its certificate data are ancillary), has an exactly certified negative rational expectation on $\sigma$. With
$g=-\tr(Y_*\sigma)$ and $h=\tr Y_*/8=1$,
the white-noise threshold $g/(g+h)$ of Lemma~\ref{lem:ppt3-dps} (Section~\ref{sec:ppt3-noise}) exceeds
$4.05\times10^{-3}$: exactly, $g>407/100000$, so $g/(g+h)>407/100407>81/20000$. The strict inequality and the
positivity of the dual matrices are checked in exact
arithmetic, without relying on solver tolerances.

\emph{Extremality.} $\sigma$ is the normalisation of $\rho(55,48,110)$, a state $\rho_{\tau,a}$ of the surface, so it is extreme by the
proofs below, which cover every $\rho_{\tau,a}$. (At $\sigma$ the criterion was also checked directly: with $D=B_\sigma X_DB_\sigma\dg$,
$B_\sigma=[\tilde u,v,w,e,o]$, the system of the next proof has exact rank $24$ over $\mathbb Q(\sqrt3)$.)
\end{proof}

\begin{proof}[Proof of the extremality of every $\rho_{\tau,a}$ in Corollary~\ref{cor:ppt3-extreme5} and of the symmetry blocks used in Proposition~\ref{prop:ppt3-moduli}]
\emph{Reduction.} An invertible local filter $F$ maps $\rho$ to $F\rho F\dg$, ranges to ranges, and each $\rho^{T_c}$ to a state
of the same kind, since $(FDF\dg)^{T_c}=\tilde FD^{T_c}\tilde F\dg$ with $\tilde F$ the filter obtained by complex-conjugating the $c$-th
factor of $F$. So $D\mapsto FDF\dg$ maps the face-direction space of $\rho$ onto that of $F\rho F\dg$, and likewise the tangent
constraints below. By Theorem~\ref{thm:ppt3-orbit}(b) it suffices to take $a=1$. Write $\rho=\rho_\tau$, $Q=\tau+\sqrt{1+\tau^2}\in(\sqrt3,1+\sqrt2)$, so that
$\tau=(Q^2-1)/(2Q)$, $\kappa=(6Q^2-Q^4-1)/(4Q^2)$,
\[
E=\tfrac14(6Q^2-Q^4-1),\quad O=\frac{6Q^2-Q^4-1}{2Q(Q^2-1)},\quad G=\tfrac12Q(5-Q^2),\quad V=\tfrac12Q(Q^2-1)(6-Q^2).
\]

\emph{Kernel of $\rho^{T_A}$.} With $\omega$ as in \eqref{eq:ppt3-vectors}, the two odd vectors $k_1,k_2$, in the basis
$(\ket{001},\ket{010},\ket{100},\ket{111})$, and the even vector $k_3$, in the basis $(\ket{000},\ket{011},\ket{101},\ket{110})$,
\begin{align*}
k_1&=\bigl((3-i\sqrt3)Q,\ (3+i\sqrt3)Q,\ 3(Q^2-1),\ 0\bigr),\\
k_2&=\bigl(Q(15-3i\sqrt3-(3+i\sqrt3)Q^2),\ Q(15+3i\sqrt3-(3-i\sqrt3)Q^2),\ 0,\ 12\bigr),\\
k_3&=\bigl(-3Q(Q^2-1),\ -6Q,\ (1-\bar\omega)Q^2+3(1+\bar\omega),\ (1-\omega)Q^2+3(1+\omega)\bigr)
\end{align*}
satisfy $\rho^{T_A}k_j=0$ identically in $Q$ and are linearly independent for $Q>1$ (the entries $3(Q^2-1)$ of $k_1$, $12$ of $k_2$ and $-6Q$ of $k_3$ are
their pivots). Since $\rank\rho^{T_A}=5$ on the window, $K_A=[k_1,k_2,k_3]$ spans $\ker\rho^{T_A}$ there. The cyclic shift $\Pi$ satisfies
$\Pi\rho\Pi\dg=\rho$ and maps cut $A$ to cut $B$, so $\ker\rho^{T_B}=\Pi\ker\rho^{T_A}$, and similarly for $C$.

\emph{Sectors.} $\Pi$, the parity $Z^{\otimes3}$ and the reflection-conjugation $K\circ S_{BC}$ fix $\rho$. Take the basis
$(u,v,w,e,o)$ of $R$ extended by $u_\perp=3\ket{000}+\ket{011}+\ket{101}+\ket{110}$, $\bar e$, $\bar o$ (spanning $R^\perp$); the matrix units $v_iv_j\dg$ have a $\Pi$-charge
$q\in\mathbb Z_3$ ($\Pi v_iv_j\dg\Pi\dg=\omega^qv_iv_j\dg$) and a parity. The real solution space of the face system (and the kernel of the
constraint map below) is invariant under conjugation by $\Pi$, so it is the sum of its charge components, and for a component of charge $q$ the conditions at
$B$ and $C$ are the one at $A$ multiplied by phases. Hence it is enough to impose the condition at the cut $A|BC$ on each sector
separately. The charges $q=1$ and $q=2$ are exchanged by $D\mapsto D\dg$.

\emph{Extremality from the $2\times4$ theorem.} Let $\rho$ be a $\rho_{\tau,a}$. Its profile is $(5,5,5,5)$ and it is GME, hence entangled across $A|BC$.
Suppose $\rho$ were not extreme in the set $\mathcal P_A$ of states with $\sigma^{T_A}\succeq0$. A ray from $\rho$ inside the face of $\mathcal P_A$ generated by
$\rho$ leaves the face at a state $\rho'$ with $\operatorname{ran}\rho'\subseteq\operatorname{ran}\rho$, $\operatorname{ran}\rho'^{T_A}\subseteq\operatorname{ran}\rho^{T_A}$
and $\min(\rank\rho',\rank\rho'^{T_A})\le4$. Such a $2\times4$ PPT state is separable~\cite{KCKL,HLVC} (if only $\rank\rho'^{T_A}\le4$, apply
the theorem to the PPT state $\rho'^{T_A}$), so $\rho'=\sum_i\ket{a_ib_i}\bra{a_ib_i}$ with
$\ket{a_1b_1}\in\operatorname{ran}\rho$ and $\ket{\bar a_1b_1}\in\operatorname{ran}\rho^{T_A}$. Put $\rho''=\rho-\epsilon\ket{a_1b_1}\bra{a_1b_1}$
with $\epsilon>0$ maximal such that $\rho''\succeq0$ and $\rho''^{T_A}\succeq0$. Then $\min(\rank\rho'',\rank\rho''^{T_A})\le4$, so $\rho''$
is separable across $A|BC$ by~\cite{KCKL,HLVC} again, and so is $\rho=\rho''+\epsilon\ket{a_1b_1}\bra{a_1b_1}$, contradicting entanglement. Hence $\rho$ is extreme in
$\mathcal P_A$, and therefore in the fully-PPT set. The second step alone shows that $\rho$ is an edge state across $A|BC$: a product vector in
$\operatorname{ran}\rho$ with partial conjugate in $\operatorname{ran}\rho^{T_A}$ would make $\rho$ separable. The same holds at every cut, and for every GME state of profile $(5,5,5,5)$.

\emph{Extremality by symmetry blocks.} The following direct computation gives a second proof and the blocks reused for Proposition~\ref{prop:ppt3-moduli}.
A Hermitian face direction is $D=BXB\dg$, $X$ a Hermitian $5\times5$ matrix ($25$ real parameters), with $D^{T_A}K_A=0$ (the
other cuts follow, and $K_c\dg D^{T_c}=0$ is the adjoint). Its sectors are:
charge $q=1$ and parity $+$ or $-$, each three complex parameters ($ue\dg,vo\dg,wo\dg$ and $uo\dg,ve\dg,we\dg$); and for $q=0$ the real parameters
$\operatorname{Re}/\operatorname{Im}$ of $(uv\dg,uw\dg,eo\dg)$ (parity $-$, three each) and, with parity $+$, $i(vw\dg-wv\dg)$ (one) and
$uu\dg,vv\dg,vw\dg+wv\dg,ww\dg,ee\dg,oo\dg$ (six). Each sector map $X\mapsto D^{T_A}K_A$ has a square submatrix whose determinant is
\begin{center}\small
\begin{tabular}{@{}ll@{}}\toprule
sector & minor \\\midrule
$q=1$, parity $+$ & $-108\,(Q^2-1)^2(Q^2+1)$ \\
$q=1$, parity $-$ & $432\,(Q^2-1)(Q^2+1)$ \\
$q=0$, parity $-$, real parts & $108\,(Q^2-1)(Q^2+2Q-1)$ \\
$q=0$, parity $-$, imaginary parts & $108\,i\,(Q^2-1)(Q^2-2Q-1)$ \\
$q=0$, parity $+$, $i(vw\dg-wv\dg)$ & $3i\,(Q^2+1)$ \\
$q=0$, parity $+$, $(uu\dg,\ vv\dg,\ vw\dg{+}wv\dg,\ ww\dg,\ ee\dg)$ & $1944\,i(\sqrt3+i)(Q^2+1)(Q^2-2Q-1)(Q^2+2Q-1)$ \\\bottomrule
\end{tabular}
\end{center}
On the window $Q^2-1$, $Q^2+1$ and $Q^2+2Q-1$ are positive and $Q^2-2Q-1<0$, so no minor vanishes. Hence the system has rank
$12+3+3+1+5=24$ (the sectors $q=1,2$ together contribute $12$, since each sector map is injective on its six real parameters), and its nullspace is spanned by $\rho$,
which lies in the last sector of six parameters. By the
supported-face criterion~\cite{GIM}, $\rho_{\tau,a}$ is extreme for every $\tau\in(1/\sqrt3,1)$ and $a>0$, as already shown.

\emph{Constraint differentials.} The tangent space of rank-five Hermitian matrices at $\rho$ consists of $D$ vanishing on
$\ker\rho\times\ker\rho$: the $55$ matrix units $v_iv_j\dg$ with $\min(i,j)\le5$ in the basis $(u,v,w,e,o;u_\perp,\bar e,\bar o)$. The differential of
the Schur-complement equation at cut $A$ is $L(D)=K_A\dg D^{T_A}K_A$, a Hermitian $3\times3$ matrix, and the same sector argument applies to the
$27$ equations of the three cuts. $L$ maps the sector of charge $q$ onto $3\times3$ matrices exactly when its two parity parts are onto:
the parity-$+$ part has the five outputs $(k_1,k_1),(k_1,k_2),(k_2,k_1),(k_2,k_2),(k_3,k_3)$ and the parity-$-$ part the four outputs
$(k_1,k_3),(k_2,k_3),(k_3,k_1),(k_3,k_2)$. Square submatrices of the sector maps have determinants
\begin{align*}
q=0,\ +:&\ -120932352\,Q(Q^2-1)^2(Q^2+1)^2(Q^4-6Q^2+1),\\
q=0,\ -:&\ -839808\,Q^2(Q^2-1)^4(Q^2-3)(Q^2+1),\\
q=1,\ +:&\ -725594112\,Q^2(Q^2-1)^4(Q^4-6Q^2+1),\\
q=1,\ -:&\ 1679616\,Q^2(Q^2-1)^3(Q^2+1)(Q^4-6Q^2+1)
\end{align*}
and $Q^4-6Q^2+1=(Q^2-2Q-1)(Q^2+2Q-1)$ vanishes only at $Q=1+\sqrt2$ and $Q=\sqrt2-1$. So on the window
every sector map is onto, the $27$ differentials are independent on the $55$-dimensional tangent space for every state of the surface, and the
zero set of the $27$ equations is a smooth manifold of dimension $55-27=28$. The factors $Q^2-3$ and $Q^4-6Q^2+1$ vanish exactly at the two
ends $\tau=1/\sqrt3$ and $\tau=1$ of the window. An independent exact computation of these ranks at $Q=2$ ($\tau=3/4$), without any symmetry reduction, is part of the same
check.
\end{proof}

\begin{proof}[Proof of Proposition~\ref{prop:ppt3-moduli}]
Let $\rho=\rho_{\tau,a}$ be any state of the surface, with $B$ and $K_c$ as above. Rank-five Hermitian matrices near $\rho$ form a smooth manifold of real
dimension $55$; its tangent space at $\rho$ is $\{ZB\dg+BZ\dg:Z\in\mathbb C^{8\times5}\}$. Near $\rho$ each $\rho'^{T_c}$ has five eigenvalues bounded away from
zero, and $\rank\rho'^{T_c}\le5$ is the vanishing of a $3\times3$ Hermitian Schur complement $S_c(\rho')$, a smooth
function whose differential at $\rho$ is $D\mapsto K_c\dg D^{T_c}K_c$, with $K_c$ a basis of $\ker\rho^{T_c}$. The
$27$ real equations $S_c=0$ have independent differentials on the tangent space, by the sector minors of the proof above, so by
the implicit function theorem their common zero set is a smooth manifold of dimension $55-27=28$, the cone over
$\mathcal M$; the trace-one condition is transversal, since $\rho$ itself is a tangent direction. Nearby points
are positive semidefinite with positive semidefinite partial transposes, since their five nonzero eigenvalues stay
positive. GME is an open condition, so it persists near $\rho$. The extremality system of
Corollary~\ref{cor:ppt3-extreme5} is built from $\operatorname{ran}\rho'$ and the kernels $\ker\rho'^{T_c}$. On $\mathcal M$ the ranks
of $\rho'$ and of every $\rho'^{T_c}$ are constant ($=5$), so these subspaces, and hence the system, vary continuously with the state.
The system has rank $24$ at $\rho$; rank is lower semicontinuous and the state itself is always a null direction, so the rank
stays $24$ nearby. The tangent to the
local-filter orbit, $\{X\rho+\rho X\dg\}$ with $X=X_1\otimes I\otimes I+I\otimes X_2\otimes I+I\otimes I\otimes X_3$,
has dimension $19$ including scale: the effective local-filter group has real dimension $24-4-1=19$
($\mathrm{GL}(2,\mathbb C)^{\times3}$ modulo the scalings $\alpha_1\alpha_2\alpha_3=1$ and the overall phase), and the
stabiliser of $\rho$ lies in the stabiliser of its range $R_{1/a}$, which is finite modulo scalars
(Corollary~\ref{cor:ppt3-chirality}(b), the group $S_3$).
The orbit tangent lies in the
cone's tangent space, since invertible local filters preserve ranks of $\rho$ and of each $\rho^{T_c}$ and so map the cone
over $\mathcal M$ to itself. The orbit dimension remains $19$
nearby: rank cannot fall, and $19$ is the maximum, so the orbits in
$\mathcal M$ have dimension $18$ and form a foliation with nine-dimensional local leaf space.
\end{proof}

\section{Exact data for one support, both phases}\label{sec:ppt3-app-coexistence}

This appendix contains the proofs of Theorem~\ref{thm:ppt3-phase-coexistence} and
Proposition~\ref{prop:ppt3-exposed-transition}, with the exact rational coefficient matrices and principal minors. The
notation, including $\Gamma_z$ of \eqref{eq:ppt3-gamma}, is that of Section~\ref{sec:ppt3-coexistence}.

In the basis $(u,v,w,e,o)$,
\begin{equation}\label{eq:ppt3-bs-same-support}
 X_{\rm BS}=\frac1{54375}
 \begin{pmatrix}
 1232&0&0&0&0\\
 0&30493&32&0&0\\
 0&32&1693&0&0\\
 0&0&0&3700&0\\
 0&0&0&0&925
 \end{pmatrix}.
\end{equation}

\begin{proof}[Proof of Theorem~\ref{thm:ppt3-phase-coexistence}]
The norms
$\|Bq_1(2)\|^2=1250$ and
$\|Bq_1(2i)\|^2=290$ give by direct expansion
\[
 \sigma_{\rm BS}
   =\tfrac13\bigl(\ket{111}\bra{111}+\Gamma_2+\Gamma_{2i}\bigr)
   =BX_{\rm BS}B^\dagger ,
\]
with $X_{\rm BS}$ as in \eqref{eq:ppt3-bs-same-support}; $\sigma_{\rm BS}$ is biseparable because each
$\Gamma_z$ is. The coefficient matrix
is positive definite and, after applying $B$, gives a trace-one state on $R$.
The state is invariant under cyclic permutations and global parity. For the unnormalised
matrix $54375\,\sigma_{\rm BS}^{T_A}$,
the leading principal minors of its even and odd
parity blocks, in the orders of
\eqref{eq:ppt3-minors}, are respectively
\[
\begin{array}{c|rrrr}
&\Delta_1&\Delta_2&\Delta_3&\Delta_4\\\hline
\mathrm{even}&1232&4558400&11751097500&13766553000000\\
\mathrm{odd}&2618&4698075&8087598750&83202897843750
\end{array}.
\]
All are strictly positive, so Sylvester's criterion
and cyclic invariance give
$\sigma_{\rm BS}^{T_c}\succ0$ for each cut.
The exact expansion and minors are checked
independently (Appendix~\ref{sec:ppt3-verifiers}).

The state $\rho_{10/3}$ has range $R$,
rank five and strictly positive partial transposes,
but is GME by \eqref{eq:ppt3-h5-value}.
For the interior assertion, let $\mu$ be a probability
measure of full support on the complex projective line $\mathbb{CP}^1$ (for example
normalised Fubini--Study measure) and
$M=\int q_1(z)q_1(z)^\dagger\,d\mu/\|Bq_1(z)\|^2$, the
coefficient matrix of a biseparable state on the
$A|BC$ product curve. If a Hermitian $H$ has
$q_1(z)^\dagger Hq_1(z)\ge0$ for all $z$ and
$\tr(HM)=0$, the nonnegative integrand vanishes on the full support of the measure, so the polynomial vanishes identically. The
five entries of $q_1(z)$ form a basis of the polynomials of degree at most four. Their products are therefore an
invertible change of basis from the independent monomials $\bar z^i z^j$, $0\le i,j\le4$, so $H=0$. Hence $M$ is
interior to the cone of supported $A|BC$-biseparable
coefficient matrices. (A measure on the unit circle
would not do: $H_5$ vanishes there.)
Mixing a sufficiently small amount
of such a moment matrix into $\sigma_{\rm BS}$
produces a biseparable state interior to the
supported biseparable cone, while retaining
strict PPT and rank five. GME is open and $\rho_{10/3}^{T_c}\succ0$,
so a relative neighbourhood of $\rho_{10/3}$
consists of fully-PPT GME states. Both regions therefore have relative interior in the same PPT spectrahedron.
\end{proof}

In the basis $(u,v,w,e,o)$,
\begin{equation}\label{eq:ppt3-xdelta}
 X_\partial=\frac1{3360}
 \begin{pmatrix}
 293&0&0&0&0\\
 0&577&-105&0&0\\
 0&-105&273&0&0\\
 0&0&0&188&0\\
 0&0&0&0&76
 \end{pmatrix}.
\end{equation}

\begin{proof}[Proof of Proposition~\ref{prop:ppt3-exposed-transition}]
For the rays in \eqref{eq:ppt3-exposed-state},
$|z|=1$ at $z=\pm1,\pm i$, while at
$z=\pm i\sqrt5$ one has $r=5$ and
$\operatorname{Re}z=0$. The factorisation
\eqref{eq:ppt3-h5} therefore vanishes on every
summand, proving $\tr(H_5X_\partial)=0$.
Their squared physical norms
$\|Bq_1(1)\|^2=56$, $\|Bq_1(i)\|^2=32$
and $\|Bq_1(i\sqrt5)\|^2=576$ give the
coefficient matrix \eqref{eq:ppt3-xdelta} upon
cyclic and sign averaging. It has
$\tr(BX_\partial B^\dagger)=1$ and
$X_\partial\succ I_5/100$ (a diagonal and
$2\times2$ block check).

For a quantitative PPT bound, the leading
principal minors of the even and odd
$4\times4$ parity blocks of
$5\,(3360\,\sigma_\partial^{T_A})-56I_8$,
in the orders of \eqref{eq:ppt3-minors}, are
\[
\begin{array}{c|rrrr}
 &\Delta_1&\Delta_2&\Delta_3&\Delta_4\\\hline
\mathrm{even}&1409&1163516&654702384&109183407616\\
\mathrm{odd}&1689&1363796&97132144&7881970176
\end{array}.
\]
They are positive. Hence
$\sigma_\partial^{T_A}\succ I/300$ by
Sylvester's criterion, and cyclic symmetry
gives the other cuts.

Let $\Delta=\sigma_{\rm BS}-\sigma_\partial$
and $X_\Delta=X_{\rm BS}-X_\partial$.
Direct rational arithmetic gives
\[
 \|\Delta\|_2^2=
 \frac{575589900887}{2318006250000}<\frac14 .
\]
Partial transposition preserves the
Hilbert--Schmidt norm, so
$\|\Delta^{T_c}\|_\infty<1/2$.
Likewise $B^\dagger B=\operatorname{diag}(4,1,3,3,3)
\succeq I_5$ implies
$\|X_\Delta\|_\infty\le\|X_\Delta\|_2
\le\|\Delta\|_2<1/2$.
For $-1/200\le t\le0$, it follows that
\[
 \sigma(t)^{T_c}\succ
   \left(\frac1{300}-\frac1{400}\right)I
      =\frac1{1200}I,\qquad
 X_\partial+tX_\Delta\succ
   \left(\frac1{100}-\frac1{400}\right)I_5 .
\]
For $0\le t\le1$ strict PPT and rank five
follow by convexity from the corresponding
properties of the two endpoints. Biseparability in that range follows from the explicit decompositions of the endpoints.
Finally, for $t<0$ the supported witness has
\[
 \tr\!\bigl(H_5(X_\partial+tX_\Delta)\bigr)
 =t\,\frac{25082}{54375}<0,
\]
which excludes biseparable mixtures.
The arithmetic is checked independently (Appendix~\ref{sec:ppt3-verifiers}).
\end{proof}

\section{Computations for the noise certificates}\label{sec:ppt3-app-noise}

This appendix contains the proofs of the statements of Section~\ref{sec:ppt3-noise} that are not given there. The rational certificate data
themselves are ancillary (Appendix~\ref{sec:ppt3-verifiers}).

\begin{proof}[Proof of Lemma~\ref{lem:ppt3-dps}]
For $x=a\otimes b$ write $a\otimes a\otimes b=Vc$. Partial transposition of the product projector gives
\[
\|a\|^2\langle x|Y|x\rangle=c\dg Pc+\langle a a\bar b|Z_1|a a\bar b\rangle+\langle a\bar a b|Z_2|a\bar a b\rangle\ \ge0 .
\]
One certificate per orbit of cuts under the symmetries of $Y$ suffices, and cyclic invariance covers all cuts on
three qubits; thus $Y$ is nonnegative on every pure state that is a product across some cut, hence on the convex hull
$\mathrm{BS}$. The noise threshold is the condition $-(1-p)g+ph<0$. The second bound follows from
$|\tr(Y(\sigma-\rho))|\le(u-l)D(\sigma,\rho)$ for $\sigma\in\mathrm{BS}$. For the last assertion, fix a cut and a unit vector
$\phi$, and set $A_\phi=(|\phi\rangle\langle\phi|)^{T_c}$.
If $\lambda_1\ge\lambda_2\ge\cdots$ are the squared
Schmidt coefficients of $\phi$ at that cut (padded by zeros if necessary), the spectrum
of $A_\phi$ consists of $\lambda_i$,
$\pm\sqrt{\lambda_i\lambda_j}$ for $i<j$, and zeros.
Since $\lambda_2\le1-\lambda_1$, its spectral width is bounded by
\[
\lambda_{\max}(A_\phi)-\lambda_{\min}(A_\phi)
=\lambda_1+\sqrt{\lambda_1\lambda_2}
\le\max_{0\le x\le1}\left[x+\sqrt{x(1-x)}\right]
=\frac{1+\sqrt2}{2}.
\]
For $\Delta=\sigma-\rho$ with $\tr\Delta=0$ and
$D(\rho,\sigma)=r$, the positive and negative parts
of $\Delta$ each have trace $r$. Consequently
$|\tr(A_\phi\Delta)|$ is at most $r$ times the
spectral width of $A_\phi$. Since $\phi$ is arbitrary,
\[
 \langle\phi|\sigma^{T_c}|\phi\rangle
 \ge c-\frac{1+\sqrt2}{2}\,r>0.
\]
The same reasoning applies to every cut.
\end{proof}

\begin{proof}[Proof of Corollary~\ref{cor:ppt3-noise}]
Apply Lemma~\ref{lem:ppt3-dps}. The exact verification checks $Z_1,Z_2,P\succ0$, cyclic symmetry, all $144$
coefficients of the product-vector identity, outward spectral bounds and every number in the table. For local noise a
check at the single rate $q_0$ suffices: $\mathcal D_{q_0}$ on all sites is obtained from any rates $q_j\le q_0$ by further
local depolarisation of rates $(q_0-q_j)/(1-q_j)$, which preserves $\mathrm{BS}$, so a GME output at $q_0$ forces a GME
input. For uniformity in $t$, the numerator of $\rho_t$ is affine in $t$, so $7(t+3)\tr(Y\rho_t)$ is affine in $t$, and so are
$7(t+3)\bigl[(1-p)\tr(Y\rho_t)+ph\bigr]$ and $7(t+3)\bigl[\tr(Y\rho_t)+r(u-l)\bigr]$ for the fixed $Y$, with $h=\tr Y/8$,
$\operatorname{spec}Y\subset[l,u]$ and fixed $p$ and ball radius $r$. An affine function that is negative at both endpoints of the
rational bracket $[3.2567,3.3886]$, which contains \eqref{eq:ppt3-interval}, is negative throughout it, and $7(t+3)>0$. By
Lemma~\ref{lem:ppt3-dps} (the noise threshold $-(1-p)g+ph<0$ and the distance bound $D(\rho,\mathrm{BS})\ge g/(u-l)$), the two
endpoint checks therefore certify the stated interval-wide white-noise and GME-ball bounds.
\end{proof}

\begin{proof}[Proof of Corollary~\ref{cor:ppt3-noise-reused}]
Write $Y=Y_*$ for the cyclic rational witness of the
certificate data. Its exact level-two certificate proves $\tr(Y\sigma)\ge0$ on
all biseparable states. The Pythagorean identity
$312^2+266^2=410^2$ makes the coefficients of
$\rho_*$ rational. The exact certificate evaluation gives
\[
  g:=-\tr(Y\rho_*)>\frac{173}{40000},\qquad
  h:=\frac{\tr Y}{8}=1,\qquad
  -\frac14 I\prec Y\prec\frac{31}{10}I .
\]
Thus the spectral width of $Y$ is less than $67/20$.
These inequalities give the white-noise and trace-distance bounds.

For (i), at $p=2153/500000$,
\[
 (1-p)\frac{173}{40000}-p
 =\frac{7531}{20000000000}>0,
\]
so $\tr[Y((1-p)\rho_*+pI/8)]<0$. The inequality
also holds for smaller $p$. Part (ii) is the
exact local-noise check at $q=81/50000$:
all smaller, unequal local rates can be followed by
additional local depolarisation to reach the checked
uniform output. Local channels preserve biseparability,
and these channels preserve full PPT.

For (iii), let $p_0=32/10000$ and $r_0=331/1000000$.
The witness margin remains strictly negative throughout
the trace-distance ball because
\[
 (1-p_0)\frac{173}{40000}-p_0
           -\frac{67}{20}r_0
  =\frac{231}{100000000}>0.
\]
Since $\rho_*^{T_c}\succeq0$, mixing with white
noise also gives $\rho_\circ^{T_c}\succeq p_0I/8$.
The PPT stability estimate in
Lemma~\ref{lem:ppt3-dps} applies because
\[
 \frac{p_0}{8}-\frac{1+\sqrt2}{2}r_0
 =\frac{469-331\sqrt2}{2000000}>0.
\]
Finally, (iv) follows directly from
\[
 \frac{173}{40000}
 -\frac{67}{20}\frac{129}{100000}
 =\frac7{2000000}>0.
\]
The rational inequalities, including the enclosure
$\sqrt2<99/70$, are checked exactly from the certified data
(Appendix~\ref{sec:ppt3-verifiers}).
\end{proof}

We now give the state $\tilde\rho$ explicitly. With $v,w,e,o$ as in \eqref{eq:ppt3-vectors}, write $\ket g=\ket{000}$ and
$\ket h=\ket{011}+\ket{101}+\ket{110}$. Then
\begin{equation}\label{eq:ppt3-tilde-rho}
\begin{aligned}
10^7\,\tilde\rho={}&955592\,\ket g\bra g-930700\,(\ket g\bra h+\ket h\bra g)+906857\,\ket h\bra h\\
&+519407\,\ket e\bra e+208476\,\ket{\bar e}\bra{\bar e}\\
&+390647\,\ket v\bra v-520181\,(\ket v\bra w+\ket w\bra v)+693325\,\ket w\bra w\\
&+556418\,\ket o\bra o+104\,\ket{\bar o}\bra{\bar o}.
\end{aligned}
\end{equation}
These ten coefficients are the state's coordinates in the invariant basis of the cyclic, parity and $K\circ S_{BC}$
symmetries; the witness data of Corollary~\ref{cor:ppt3-robust} are ancillary.

\begin{proof}[Proof of Corollary~\ref{cor:ppt3-robust}]
Lemma~\ref{lem:ppt3-dps} applies. The exact verification checks trace one, full rank, cyclic and parity invariance, all
three partial-transpose eigenvalue lower bounds, and the witness and tolerance inequalities. The PPT-only ball radii here
and in the table of Corollary~\ref{cor:ppt3-noise} follow from those eigenvalue bounds and the stability estimate of
Lemma~\ref{lem:ppt3-dps}; their rational comparisons are checked against the same constants.
\end{proof}

\begin{proof}[Proof of Proposition~\ref{prop:ppt3-noise-sharp}]
The lower bounds are Corollaries~\ref{cor:ppt3-noise-reused}, \ref{cor:ppt3-noise} and~\ref{cor:ppt3-robust}. For the
upper bounds, let $\mathcal C(X)=\frac13\sum_{k=0}^2\Pi^kX\Pi^{-k}$ average over the cyclic qubit permutation $\Pi$. The
certificate data give, for each of the three states, three orbit representatives: a Gaussian-rational qubit vector $a_r$ and a
positive definite Gaussian-rational $4\times4$ matrix $\Lambda_r$. The states are invariant under global parity $Z^{\otimes3}$ and
the antiunitary $K\circ S_{BC}$, which preserve the cut $A|BC$; applying them to $|a_r\rangle\langle a_r|\otimes\Lambda_r$ gives at most four terms
per representative, and $\theta$ is the sum of all of them (at most twelve terms). Thus $\theta=\sum_i|a_i\rangle\langle a_i|\otimes\Lambda_i$ with $\Lambda_i\succ0$
is separable across $A|BC$.
At $p$ equal to the stated upper bound, the residual $\Delta=(1-p)\rho+pI/8-\mathcal C(\theta)$ is cyclically invariant and, with $c=\tr\Delta/8>0$,
satisfies $\|\Delta-cI\|_2\le c$ exactly. By the bipartite Gurvits--Barnum bound~\cite{GB}, this makes $\Delta$
separable across $A|BC$. Hence
$(1-p)\rho+pI/8=\mathcal C(\theta+\Delta)$ is an equal mixture of states separable across the three cuts, so it is biseparable at
$p$ equal to the stated upper bound. For every larger $p'$, $(1-p')\rho+p'I/8$ is a convex mixture of the state at $p$ with the
separable state $I/8$, hence biseparable too.
\end{proof}

\section{The quartic moment cone of a minimal-rank support}\label{sec:ppt3-app-moment}

We first rule out cut-product vectors in the kernel. The remaining three kernel constraints then force a unique
product ray for each qubit direction, producing the quartic curves of Theorem~\ref{thm:ppt3-universal-moments}.

\begin{lemma}[product kernel]\label{lem:ppt3-product-kernel}
\cls{proved}
A PPT operator on $\mathbb C^2\otimes\mathbb C^4$ whose kernel contains a nonzero product vector is separable. Hence the
kernel of a fully-PPT three-qubit state that is not separable across a cut contains no product vector across that cut,
and the range of every rank-five fully-PPT GME state satisfies the hypothesis of Theorem~\ref{thm:ppt3-universal-moments}.
\end{lemma}

\begin{proof}
This is a standard reduction, cf.~\cite{KCKL}; we give a short proof. Separability and PPT are invariant under local unitaries, so we may take the kernel vector to be $\ket0\ket3$. Write
$\rho=\begin{pmatrix}A&B\\B\dg&C\end{pmatrix}$ in qubit blocks. Then $\rho\ket{0,3}=0$ gives $A\ket3=0$ and $B\dg\ket3=0$.
Since $\langle0,3|\rho^{T_A}|0,3\rangle=\langle0,3|\rho|0,3\rangle=0$ and $\rho^{T_A}=\begin{pmatrix}A&B\dg\\B&C
\end{pmatrix}\succeq0$, also $\rho^{T_A}\ket{0,3}=0$, so $B\ket3=0$. Hence $\rho\ket{1,3}=\rho^{T_A}\ket{1,3}=
\ket1\otimes C\ket3$. If $C_{33}=\langle3|C|3\rangle=0$ then $C\ket3=0$ and $\rho$ lives on $\mathbb C^2\otimes\mathbb C^3$.
Otherwise put $u=C\ket3/\sqrt{C_{33}}$ and $\rho'=\rho-\ket1\bra1\otimes\ket u\bra u$. For any $M\succeq0$ and vector $f$
with $\langle f|M|f\rangle>0$, the Schur complement $M-M\ket f\bra fM/\langle f|M|f\rangle$ is positive semidefinite. With
$f=\ket{1,3}$ this applies to $M=\rho$ and to $M=\rho^{T_A}$, and both subtracted terms equal $\ket1\bra1\otimes\ket u\bra u$,
which is invariant under partial transposition. So $\rho'$ is PPT, and $\rho'\ket{j,3}=0$ for $j=0,1$: $\rho'$ lives on
$\mathbb C^2\otimes\mathbb C^3$. PPT operators on $\mathbb C^2\otimes\mathbb C^3$ are separable~\cite{HHH96}, so $\rho=\rho'+
\ket1\bra1\otimes\ket u\bra u$ is separable. For the consequence, a rank-five fully-PPT state with a product vector in
its kernel at some cut is separable across that cut, hence biseparable.
\end{proof}

\begin{proof}[Proof of Theorem~\ref{thm:ppt3-universal-moments}]
Fix a cut; after relabelling the qubits it is $A|BC$. Let $N$ be an $8\times3$ basis matrix of $R^\perp$. A vector
$a\otimes b$ lies in $R$ if and only if $M(a)b=0$, where $M(a)=N\dg(a\otimes I_4)$ is a $3\times4$ matrix depending
linearly on $a$. If $\rank M(a)<3$ for some $a\ne0$, a nonzero $y$ has $y\dg M(a)=0$, so $k=Ny\in R^\perp$ is orthogonal
to $a\otimes\mathbb C^4$, that is $k=a^\perp\otimes b'$ is a product vector, contrary to the hypothesis. So the pencil
$M(a)=a_0M_0+a_1M_1$ has rank three for every $a\ne0$.

By the Kronecker canonical form of matrix pencils~\cite[Ch.~XII]{Gant}, there are invertible $P$ and $Q$ with
$PM(a)Q$ block diagonal, with blocks $L_\varepsilon$ ($\varepsilon\times(\varepsilon+1)$, with entries $a_0$ on the
diagonal and $a_1$ on the superdiagonal), transposes $L_\eta^T$, and a regular part. Full row rank excludes every
$L_\eta^T$ block. A regular block of positive size has a homogeneous determinant that vanishes somewhere on
$\mathbb{CP}^1$, contradicting constant rank. Each remaining $L$ block contributes one kernel dimension;
nullity one therefore forces the single block $L_3$. Its kernel is spanned by
$(a_1^3,-a_0a_1^2,a_0^2a_1,-a_0^3)^T$, so $\ker M(a)$ is spanned by $b(a)=Q\,\kappa(a)$ with $\kappa(a)$ this cubic vector.
The product rays in $R$ are therefore the rays of $a\otimes b(a)=(I_2\otimes Q)(a\otimes\kappa(a))$. The eight entries of
$a\otimes\kappa(a)$ are $\pm$ quartic monomials and every one of the five monomials occurs, so $a\otimes b(a)=L\nu(a)$ for
an $8\times5$ matrix $L$ of rank five. Its column space is spanned by the curve, lies in $R$ and has dimension five, so it
equals $R$, and $L=BT_A$ for an invertible $T_A$. (This is the case of a maximal completely entangled kernel
in~\cite[Prop.~24]{CD}.) The other cuts are the same with the qubits relabelled.

A positive decomposition of an operator supported on $R$ uses only vectors in $R$, since each positive summand has zero
expectation on $R^\perp$. A biseparable $BXB\dg$ is therefore a positive combination of cut-$c$ product projectors in
$R$, that is of $BT_c\nu(a)\nu(a)\dg T_c\dg B\dg$, which gives~\eqref{eq:ppt3-universal-bs} (the converse direction is
immediate). The trace-one product projectors on $R$ form a compact set, so their cone is closed, and Carath\'eodory's
theorem in the $25$-dimensional real space of Hermitian $5\times5$ matrices bounds the number of terms. The dual
statement~\eqref{eq:ppt3-universal-dual} is the definition of the dual of a sum of cones generated by the rank-one
matrices $T_c\nu(a)\nu(a)\dg T_c\dg$.
\end{proof}

\section{Supported witnesses extend by a kernel penalty}\label{sec:ppt3-kernel}

A supported witness need not control product vectors outside its support. Choose $K$ with kernel equal to the support.
Adding a positive multiple of $K\dg K$ penalises their component outside $\ker K$ while leaving the target state's expectation unchanged. A small identity shift makes
positivity strict on the supported product rays, allowing compactness to control the rest.

\begin{lemma}[kernel-penalty witness lift]\label{lem:ppt3-kernel-penalty}
\cls{proved}
Let $C$ be any compact set of unit vectors in a finite-dimensional Hilbert space, let $K$ be a matrix (below, $\|\cdot\|_\infty$ is the operator norm and $\|\cdot\|$ the Euclidean vector norm) and
$\mathcal R=\ker K$. If a Hermitian $W_0$ is strictly positive on $C\cap\mathcal R$, then
\[
W=W_0+\lambda K\dg K
\]
is nonnegative on all of $C$ for some finite $\lambda\ge0$. No constant-rank or smoothness hypothesis is needed. More
generally, if $W_0$ is merely nonnegative on $C\cap\mathcal R$ and $\tr(W_0\rho)=-g<0$ for a density matrix $\rho$ supported
on $\mathcal R$, then $W_0+\epsilon I+\lambda K\dg K$ is a global witness for the free set $C$ that detects $\rho$, for every
$0<\epsilon<g$ and some finite $\lambda$.
\end{lemma}
\begin{proof}
On $C$ put $f(v)=v\dg W_0v$ and $h(v)=\|Kv\|^2$. The compact set $\mathcal B=\{v\in C:f(v)\le0\}$ is disjoint from
$\{h=0\}$. If $\mathcal B$ is nonempty, $\eta=\min_{\mathcal B}h>0$, so $\lambda\ge\|W_0\|_\infty/\eta$ makes $f+\lambda h\ge0$ there;
outside $\mathcal B$ it is already positive. If $\mathcal B$ is empty take $\lambda=0$. For the nonstrict version apply
the argument to $W_0+\epsilon I$. The kernel penalty has zero expectation on $\rho$, leaving $-g+\epsilon<0$.
\end{proof}

For biseparability, $C$ is the union of the normalised cut-product vectors over the three cuts.
The lemma also applies to other free pure-state sets, such as fully product vectors or stabiliser vectors.

\begin{proposition}[quantitative form]\label{prop:ppt3-kernel-quant}
\cls{proved}
Let $C$ be the set of normalised cut-product vectors $U_c^\dagger(a\otimes b)$, $c\in\{A,B,C\}$, $a$ a unit
one-qubit vector, where $U_c$ moves qubit $c$ to the first tensor factor. Put
$M_c(a)=KU_c^\dagger(a\otimes I)$ and $\Omega_c(a)=(a^\dagger\otimes I)U_cW_0U_c^\dagger(a\otimes I)$. Suppose
$\|W_0\|_\infty\le L$, $\Omega_c(a)\succeq\alpha I$ on $\ker M_c(a)$, and every nonzero eigenvalue of $M_c(a)\dg M_c(a)$ is at least
$\mu$, uniformly in $a$ and $c$, with constants $L,\alpha,\mu>0$. Then
\begin{equation}\label{eq:ppt3-lambda}
\lambda=\frac{L+L^2/\alpha}{\mu}
\end{equation}
suffices in Lemma~\ref{lem:ppt3-kernel-penalty}. If additionally $\|K\|_\infty^2\le\beta$ and $\tr(W_0\rho)=-g<0$ for $\rho$
supported on $\ker K$, then, for the lifted witness $W=W_0+\lambda K\dg K$,
\begin{equation}\label{eq:ppt3-dist-lift}
D(\rho,\mathrm{BS})\ge\frac g{2L+\lambda\beta}.
\end{equation}
\end{proposition}
\begin{proof}
Split $b=u+v$ into its components in $\ker M_c(a)$ and its orthogonal complement. The conditional quadratic form
$(a\otimes b)\dg U_c(W_0+\lambda K\dg K)U_c\dg(a\otimes b)$ is bounded below by
\[
\alpha\|u\|^2-2L\|u\|\|v\|+(\lambda\mu-L)\|v\|^2\ \ge0
\]
when $\lambda\mu-L\ge L^2/\alpha$, which is \eqref{eq:ppt3-lambda}. For \eqref{eq:ppt3-dist-lift}, the spectrum of $W$
lies in $[-L,L+\lambda\beta]$, $\tr(W\rho)=-g$ and $\tr(W\sigma)\ge0$ on $\mathrm{BS}$; the distance bound of
\eqref{eq:ppt3-stability}, whose proof uses only these facts, gives \eqref{eq:ppt3-dist-lift}. The same argument works for any finite union of
product cuts.
\end{proof}

Compactness supplies positive $\alpha$ and $\mu$ for constant-rank pencils with strict supported positivity; the next
remark gives explicit bounds for the family.

\begin{remark}[hand-checkable constants]\label{rem:ppt3-hand-radius}
\cls{proved; CF-5}
Without any solver output, Proposition~\ref{prop:ppt3-kernel-quant} gives at $t=10/3$
\begin{equation}\label{eq:ppt3-hand}
D(\rho_{10/3},\mathrm{BS})>\frac1{20\,000\,000},\qquad
\rho_{10/3}^{T_c}\succ I/2000,
\end{equation}
so every physical state in that closed trace-distance ball is fully PPT and GME. Corollary~\ref{cor:ppt3-noise} is stronger
by a factor of about $5000$; \eqref{eq:ppt3-hand} uses no semidefinite
program, only explicit constants and exact minors.

\emph{Proof.} Let $T=B\dg B$ and $W_0=BT^{-1}H_5T^{-1}B\dg$. Its operator norm is $14/3$ and, by
\eqref{eq:ppt3-h5-value}, its expectation on $\rho_{10/3}$ is $-3/133$. Normalise the three rows of the annihilator $K$ of
$R$ (Appendix~\ref{sec:ppt3-app-family}, proof of Proposition~\ref{prop:ppt3-witness}) by $\sqrt3$, $\sqrt3$, $\sqrt{12}$, which gives $KK\dg=I$. For any unit $a$ the
conditional matrix $M=K(a\otimes I)$ satisfies $MM\dg\preceq I$ and $\det(MM\dg)\ge1/81$: write $|a_0|^2=p$ and
$a_0\bar a_1=z$; the minimum over $\arg z$ is
\[
\det(MM\dg)\ge\frac{1-7p+24p^2-12p^3}{36}\ge\frac1{81},
\]
because the numerator minus $4/9$ is $(5-3p)(6p-1)^2/9$. Consequently every nonzero squared singular value is at least
$1/81$. Take $U_c$ to be the power of the cyclic qubit permutation that moves $c$ to the first factor. The range $R$, hence
the row space of $K$, and $W_0$ are cyclically invariant, so $KU_c\dg=G_cK$ with $G_c$ unitary and $U_cW_0U_c\dg=W_0$; the
cuts $B$ and $C$ therefore give the same constants as $A$. Apply Proposition~\ref{prop:ppt3-kernel-quant} to $W_0+I/100$ with $L=5$, $\alpha=1/100$, $\mu=1/81$ and
$\beta=1$, so $\lambda=202905$. The global witness $W=W_0+I/100+202905\,K\dg K$ has expectation $-167/13300$ on
$\rho_{10/3}$ and spectral range at most $202915$, which gives the strict distance bound in \eqref{eq:ppt3-hand}. The
partial-transpose bound is Sylvester's criterion for the two blocks of \eqref{eq:ppt3-minors} after subtracting $I/2000$,
whose leading minors are positive rationals. The ball statement then follows from the last part of
Lemma~\ref{lem:ppt3-dps}, since $1/20\,000\,000<2/(2000(1+\sqrt2))$.
\end{remark}

\section{Proof dependencies and exact certificate data}\label{sec:ppt3-verifiers}

The existence and PPT-window theorems are analytic: the displayed formulas suffice to check the product curves,
square factorisations, chirality and sign bounds by hand. Extremality, the moduli count and quantitative noise bounds
also use exact computational certificates. We list their dependencies below; the preceding appendices contain the proofs.

\begin{description}
\item[Mixed $W$ class.]
Theorem~\ref{thm:ppt3-w-class} is analytic. The finite cyclic/parity decompositions, zero-hyperdeterminant
identities, surface coefficient budgets, phase-coexistence line, and full-rank state $\tilde\rho$ are independently checked in exact arithmetic. The check includes both surface branches,
the phase-coexistence line, $\tilde\rho$, nontrivial local filters, and a GHZ negative control; no numerical convex-roof minimisation is used.

\item[Extremality and moduli.]
The extremality of Corollary~\ref{cor:ppt3-extreme6} rests on an exact $24$-of-$25$ rank calculation. For every state of the
$(5,5,5,5)$ surface, the extremality of Corollary~\ref{cor:ppt3-extreme5} and the $27$ independent constraint differentials of
Proposition~\ref{prop:ppt3-moduli} rest on symmetry-block minors that are explicit polynomials in $Q$, checked symbolically and by an exact rank
computation at $Q=2$; the local-filter orbit has dimension $19$ because the stabiliser of the range is finite
(Corollary~\ref{cor:ppt3-chirality}(b)). The proofs are in Appendices~\ref{sec:ppt3-app-family} and~\ref{sec:ppt3-app-surface}. These results are computer-assisted.

\item[One support, both phases.]
Theorem~\ref{thm:ppt3-phase-coexistence} and Proposition~\ref{prop:ppt3-exposed-transition} rest on exact rational
expansions of finite cut-product mixtures, principal minors of the partial-transpose parity blocks, witness values and a
Hilbert--Schmidt bound (Appendix~\ref{sec:ppt3-app-coexistence}).

\item[Noise bounds.]
The level-two certificates use rational matrices $Y$, $Z_1$, $Z_2$ given in symmetry-adapted coordinates (cyclic symmetry, global
parity and the antiunitary $K\circ S_{BC}$): $Y$ has $10$ real parameters and each $Z_i$ is a pair of real symmetric $8\times8$
blocks ($72$ parameters). The required positivity conditions
and polynomial identities are verified in exact arithmetic. The certificate data consist of the witnesses for the family, the state $\tilde\rho$ with
its witness, and the witness $Y_*$ for the surface state. Every white-noise bound, radius and table entry of
Section~\ref{sec:ppt3-noise} (proofs in Appendix~\ref{sec:ppt3-app-noise}) is checked from these data, not inferred from solver output; this includes the spectral
enclosure $-I/4\prec Y_*\prec31I/10$, the rational value of $\tr(Y_*\sigma)$, the partial-transpose eigenvalue
bounds and the use of $\sqrt2<99/70$ in the PPT radii. Negative controls lie just beyond each stated
bound.

\item[Sharp white-noise thresholds.]
Proposition~\ref{prop:ppt3-noise-sharp} uses explicit biseparable decompositions, stored as three orbit representatives per state
and expanded under parity and $K\circ S_{BC}$. The exact checks cover positive
definiteness of every $\Lambda_r$, cyclic invariance of the residual and the Gurvits--Barnum inequality, with a negative
control at the certified GME noise level.
\end{description}

Corollary~\ref{cor:ppt3-symmetrisation}, Proposition~\ref{prop:ppt3-noise-extremal}, and the proofs of Appendices~\ref{sec:ppt3-app-moment} and~\ref{sec:ppt3-kernel} are analytic; the constants of Remark~\ref{rem:ppt3-hand-radius} are checked exactly. For publication, the
certificate data and the exact verifiers must accompany the paper as ancillary material or a version-pinned public
archive; the PDF alone does not encode the certificate matrices. Witness-search routines supply candidates; only exact checks certify the results.



\begin{thebibliography}{99}
\small\setlength{\itemsep}{0pt}
\bibitem{SB} G.~Singh and M.~Banik, On the structure of low-rank PPT entangled states violating Peres conjecture,
\href{https://arxiv.org/abs/2610.19905}{arXiv:2610.19905} (14 October 2026).
\bibitem{JMG} B.~Jungnitsch, T.~Moroder and O.~G\"uhne, Taming multiparticle entanglement, Phys.\ Rev.\ Lett.\
\textbf{106}, 190502 (2011), \href{https://arxiv.org/abs/1010.6049}{arXiv:1010.6049}.
\bibitem{NMG} L.~Novo, T.~Moroder and O.~G\"uhne, Genuine multiparticle entanglement of permutationally invariant states,
Phys.\ Rev.\ A \textbf{88}, 012305 (2013), \href{https://arxiv.org/abs/1302.4100}{arXiv:1302.4100}.
\bibitem{HK} K.-C.~Ha and S.-H.~Kye, Construction of three-qubit genuine entanglement with bipartite positive partial
transposes, Phys.\ Rev.\ A \textbf{93}, 032315 (2016), \href{https://arxiv.org/abs/1512.04693}{arXiv:1512.04693}.
\bibitem{HaKye15} K.-C.~Ha and S.-H.~Kye, Multi-partite separable states with unique decompositions and construction of
three qubit entanglement with positive partial transpose, J.\ Phys.\ A \textbf{48}, 045303 (2015),
\href{https://arxiv.org/abs/1402.5813}{arXiv:1402.5813}.
\bibitem{HanK} K.~H.~Han and S.-H.~Kye, Construction of multi-qubit optimal genuine entanglement witnesses, J.\ Phys.\ A
\textbf{49}, 175303 (2016), \href{https://arxiv.org/abs/1510.03620}{arXiv:1510.03620}.
\bibitem{ABLS} A.~Ac\'in, D.~Bru\ss, M.~Lewenstein and A.~Sanpera, Classification of mixed three-qubit states, Phys.\ Rev.\
Lett.\ \textbf{87}, 040401 (2001), \href{https://arxiv.org/abs/quant-ph/0103025}{arXiv:quant-ph/0103025}.
\bibitem{DVC} W.~D\"ur, G.~Vidal and J.~I.~Cirac, Three qubits can be entangled in two inequivalent ways, Phys.\ Rev.\ A
\textbf{62}, 062314 (2000), \href{https://arxiv.org/abs/quant-ph/0005115}{arXiv:quant-ph/0005115}.
\bibitem{CKW} V.~Coffman, J.~Kundu and W.~K.~Wootters, Distributed entanglement, Phys.\ Rev.\ A \textbf{61}, 052306 (2000),
\href{https://arxiv.org/abs/quant-ph/9907047}{arXiv:quant-ph/9907047}.
\bibitem{HLVC} P.~Horodecki, M.~Lewenstein, G.~Vidal and I.~Cirac, Operational criterion and constructive checks for the separability of low-rank density matrices, Phys.\ Rev.\ A \textbf{62}, 032310 (2000), \href{https://arxiv.org/abs/quant-ph/0002089}{arXiv:quant-ph/0002089}.
\bibitem{KL} S.~Karnas and M.~Lewenstein, Separability and entanglement in $2\times2\times N$ composite quantum systems,
Phys.\ Rev.\ A \textbf{64}, 042313 (2001), \href{https://arxiv.org/abs/quant-ph/0102115}{arXiv:quant-ph/0102115}.
\bibitem{KPBS} H.~Kampermann, D.~Bru\ss, X.~Peng and D.~Suter, Experimental generation of pseudo-bound-entanglement, Phys.\
Rev.\ A \textbf{81}, 040304(R) (2010), \href{https://arxiv.org/abs/0909.2743}{arXiv:0909.2743}.
\bibitem{TV} G.~T\'oth and T.~V\'ertesi, Quantum states with a positive partial transpose are useful for metrology, Phys.\
Rev.\ Lett.\ \textbf{120}, 020506 (2018), \href{https://arxiv.org/abs/1709.03995}{arXiv:1709.03995}.
\bibitem{Peres} A.~Peres, Separability criterion for density matrices, Phys.\ Rev.\ Lett.\ \textbf{77}, 1413 (1996),
\href{https://arxiv.org/abs/quant-ph/9604005}{arXiv:quant-ph/9604005}.
\bibitem{HHH96} M.~Horodecki, P.~Horodecki and R.~Horodecki, Separability of mixed states: necessary and sufficient
conditions, Phys.\ Lett.\ A \textbf{223}, 1 (1996), \href{https://arxiv.org/abs/quant-ph/9605038}{arXiv:quant-ph/9605038}.
\bibitem{HHH98} M.~Horodecki, P.~Horodecki and R.~Horodecki, Mixed-state entanglement and distillation: is there a
``bound'' entanglement in nature?, Phys.\ Rev.\ Lett.\ \textbf{80}, 5239 (1998),
\href{https://arxiv.org/abs/quant-ph/9801069}{arXiv:quant-ph/9801069}.
\bibitem{HHH09} R.~Horodecki, P.~Horodecki, M.~Horodecki and K.~Horodecki, Quantum entanglement, Rev.\ Mod.\ Phys.\
\textbf{81}, 865 (2009), \href{https://arxiv.org/abs/quant-ph/0702225}{arXiv:quant-ph/0702225}.
\bibitem{GT09} O.~G\"uhne and G.~T\'oth, Entanglement detection, Phys.\ Rep.\ \textbf{474}, 1 (2009),
\href{https://arxiv.org/abs/0811.2803}{arXiv:0811.2803}.
\bibitem{KRB} H.~Krishnan S.~V., A.~Ranjan and M.~Banik, State space structure of tripartite quantum systems,
Phys.\ Rev.\ A \textbf{104}, 022437 (2021), \href{https://arxiv.org/abs/2104.06938}{arXiv:2104.06938}.
\bibitem{CD} L.~Chen and D.~\v{Z}.~\DJ{}okovi\'c, Qubit-qudit states with positive partial transpose, Phys.\ Rev.\ A \textbf{86}, 062332 (2012), \href{https://arxiv.org/abs/1210.0111}{arXiv:1210.0111}.
\bibitem{Gant} F.~R.~Gantmacher, \emph{The Theory of Matrices}, Vol.~II (Chelsea, New York, 1959), Ch.~XII.
\bibitem{KCKL} B.~Kraus, J.~I.~Cirac, S.~Karnas and M.~Lewenstein, Separability in $2\times N$ composite quantum systems,
Phys.\ Rev.\ A \textbf{61}, 062302 (2000), \href{https://arxiv.org/abs/quant-ph/9912010}{arXiv:quant-ph/9912010}.
\bibitem{GB} L.~Gurvits and H.~Barnum, Largest separable balls around the maximally mixed bipartite quantum state,
Phys.\ Rev.\ A \textbf{66}, 062311 (2002), \href{https://arxiv.org/abs/quant-ph/0204159}{arXiv:quant-ph/0204159}.
\bibitem{DPS02} A.~C.~Doherty, P.~A.~Parrilo and F.~M.~Spedalieri, Distinguishing separable and entangled states, Phys.\
Rev.\ Lett.\ \textbf{88}, 187904 (2002), \href{https://arxiv.org/abs/quant-ph/0112007}{arXiv:quant-ph/0112007}.
\bibitem{DPS} A.~C.~Doherty, P.~A.~Parrilo and F.~M.~Spedalieri, Complete family of separability criteria, Phys.\ Rev.\ A
\textbf{69}, 022308 (2004), \href{https://arxiv.org/abs/quant-ph/0308032}{arXiv:quant-ph/0308032}.
\bibitem{DPS05} A.~C.~Doherty, P.~A.~Parrilo and F.~M.~Spedalieri, Detecting multipartite entanglement, Phys.\ Rev.\ A
\textbf{71}, 032333 (2005), \href{https://arxiv.org/abs/quant-ph/0407143}{arXiv:quant-ph/0407143}.
\bibitem{TG} G.~T\'oth and O.~G\"uhne, Entanglement and permutational symmetry, Phys.\ Rev.\ Lett.\ \textbf{102}, 170503
(2009), \href{https://arxiv.org/abs/0812.4453}{arXiv:0812.4453}.
\bibitem{TAHKSL} J.~Tura, R.~Augusiak, P.~Hyllus, M.~Ku\'s, J.~Samsonowicz and M.~Lewenstein, Four-qubit entangled symmetric
states with positive partial transpositions, Phys.\ Rev.\ A \textbf{85}, 060302(R) (2012),
\href{https://arxiv.org/abs/1203.3711}{arXiv:1203.3711}.
\bibitem{ATSL} R.~Augusiak, J.~Tura, J.~Samsonowicz and M.~Lewenstein, Entangled symmetric states of $N$ qubits with all
positive partial transpositions, Phys.\ Rev.\ A \textbf{86}, 042316 (2012),
\href{https://arxiv.org/abs/1206.3088}{arXiv:1206.3088}.
\bibitem{HS} M.~Huber and R.~Sengupta, Witnessing genuine multipartite entanglement with positive maps, Phys.\ Rev.\ Lett.\
\textbf{113}, 100501 (2014), \href{https://arxiv.org/abs/1404.7449}{arXiv:1404.7449}.
\bibitem{PM} M.~Piani and C.~E.~Mora, Class of positive-partial-transpose bound entangled states associated with almost any
set of pure entangled states, Phys.\ Rev.\ A \textbf{75}, 012305 (2007), \href{https://arxiv.org/abs/quant-ph/0607061}{arXiv:quant-ph/0607061}.
\bibitem{SBZL} F.~Shi, G.~Bai, X.~Zhang, Q.~Zhao and L.~Li, A minimum-cardinality genuinely unextendible product basis in
three qutrits, \href{https://arxiv.org/abs/2608.12785}{arXiv:2608.12785} (13 August 2026).
\bibitem{GIM} {\O}.~S.~Garberg, B.~Irgens and J.~Myrheim, Extremal states of positive partial transpose in a system of
three qubits, Phys.\ Rev.\ A \textbf{87}, 032302 (2013), \href{https://arxiv.org/abs/1212.5036}{arXiv:1212.5036}.
\bibitem{Companion} A.~Lewis, Companion repository (paper, verification scripts with captured output, Lean~4 formalisation),
\href{https://github.com/cmbant/ppt-gme-three-qubit}{\texttt{github.com/cmbant/ppt-gme-three-qubit}},
archived at Zenodo, \href{https://doi.org/10.5281/zenodo.24113583}{doi:10.5281/zenodo.24113583}.
\end{thebibliography}
\end{document}
